% Discrete or Continuous? — LaTeX source (Lean 4 / Mathlib companion library: substrate-undecidability-lean)
% Companion formal artifact: Substrate/*.lean (module Substrate, repository substrate-undecidability-lean)

\documentclass[11pt]{article}
\input{nova_zenodo_doi_placeholder}
\usepackage[utf8]{inputenc}
\usepackage[T1]{fontenc}
\usepackage[margin=1in]{geometry}
\usepackage{amsmath,amssymb,amsthm}
\usepackage{mathtools}
\usepackage[colorlinks=true,citecolor=blue,linkcolor=blue,urlcolor=blue]{hyperref}
\usepackage{enumitem}
\usepackage{array}
\usepackage{booktabs}
\usepackage{longtable}
\usepackage{xcolor}
\usepackage{stmaryrd}
\usepackage{tcolorbox}
\usepackage{natbib}
\usepackage{placeins}
\usepackage{cleveref}

\newtcolorbox{keybox}[1][]{
  colback=blue!4, colframe=blue!35,
  arc=3pt, boxrule=0.5pt,
  fonttitle=\bfseries, #1
}

% Allow line breaks in long Lean identifiers (namespaces, underscores, CamelCase).
\ExplSyntaxOn
\tl_new:N \l__su_leanref_tl
\NewDocumentCommand \leanref { m }
  {
    \tl_set:Nn \l__su_leanref_tl { #1 }
    \regex_replace_all:nnN { \c{_} } { \c{_} \c{allowbreak} } \l__su_leanref_tl
    \regex_replace_all:nnN { \. } { . \c{allowbreak} } \l__su_leanref_tl
    \regex_replace_all:nnN { / } { / \c{allowbreak} } \l__su_leanref_tl
    \regex_replace_all:nnN { ([a-z])([A-Z]) } { \1 \c{allowbreak} \2 } \l__su_leanref_tl
    \regex_replace_all:nnN { ([0-9])([A-Z]) } { \1 \c{allowbreak} \2 } \l__su_leanref_tl
    \texttt{\small \tl_use:N \l__su_leanref_tl }
  }
\ExplSyntaxOff

% Long \leanref tokens and inline monospace identifiers can overflow a justified
% line; relax line-breaking globally rather than hand-tune every occurrence.
\sloppy
\setlength{\emergencystretch}{3em}

% Theorem, lemma, and corollary each get their own independent counter (not shared,
% not reset per section) so that the displayed number matches this paper's own
% cross-references (``Theorem 4'', ``Lemma 2.1'', ``Corollary 4.1'') exactly, with no
% second, conflicting auto-generated number. Definitions and remarks keep the original
% per-section shared-counter behavior, since nothing elsewhere hand-references their
% numbers and that numbering was never in conflict.
\newtheorem{theorem}{Theorem}
\newtheorem{lemma}{Lemma}
\newtheorem{proposition}{Proposition}
\newtheorem{corollary}{Corollary}
\theoremstyle{definition}
\newtheorem{definition}{Definition}[section]
\newtheorem{remark}[definition]{Remark}
\theoremstyle{plain}
\renewcommand{\thelemma}{2.\arabic{lemma}}
\renewcommand{\thecorollary}{\arabic{theorem}.\arabic{corollary}}

% The theorem and corollary counters are manually reset/decremented at several points
% below (via \setcounter/\addtocounter) so that displayed numbers match this paper's
% own hand-written cross-references. hyperref's PDF-internal anchor names are derived
% from the same raw counter values by default, so a manual reset that revisits a
% previously-used counter value produces two different environments claiming the same
% PDF destination name (silently sending some \cref hyperlinks to the wrong target,
% even though the on-page displayed number is unaffected). Every colliding pair here
% falls on a different page, so prefixing the hyperref anchor with the page number
% disambiguates them completely; this is hyperref's own documented remedy for counters
% reset after they have already been used for a destination.
%
% Confirmed still necessary (2026-09-13): removing these two lines and recompiling
% 5 passes reproduces "destination with the same identifier (name{corollary.1})"
% on every pass, from the raw `corollary` counter hitting 1 three times (Corollary
% 1.1 after Theorem 1, Corollary 4.1 after the first \setcounter{corollary}{0},
% Corollary 5.1 after the second) even though \thecorollary's displayed text (which
% also folds in \arabic{theorem}) correctly disambiguates all three. Do not remove.
\renewcommand{\theHtheorem}{\thepage.\arabic{theorem}}
\renewcommand{\theHcorollary}{\thepage.\arabic{theorem}.\arabic{corollary}}

% Theorem 4' is a genuinely distinct result from Theorem 4 (a strictly sharper reduction, not
% a restatement), so it is not folded into Theorem 4 as a labeled part; but it is also not the
% *next* item in the main sequence, since Theorem 5 already occupies that slot in this paper's
% own hand-written cross-references throughout. \newtheorem* (unnumbered, from amsthm) gives it
% a real theorem environment with its own heading text fixed to "Theorem 4'" directly, so it
% never touches the shared theorem counter and cannot collide with any other numbered theorem's
% hyperref anchor.
\newtheorem*{thmfourprime}{Theorem~4\('\)}

\newcommand{\N}{\mathbb{N}}
\newcommand{\R}{\mathbb{R}}
\newcommand{\Baire}{\N \to \N}
\newcommand{\A}{\textsc{a\_lean}}
\newcommand{\AMDL}{\textsc{a\_mdl}}
\newcommand{\D}{\mathrm{D}}

\title{Discrete or Continuous?\\[0.2em]
  {\large Why No Internal Observer Can Decide, and What Can Adjudicate Instead}}
\author{Nova Spivack\\
\url{https://www.novaspivack.com}}
\date{2026}

\begin{document}
\maketitle
\NovaZenodoAfterTitle

% ============================================================================
% ABSTRACT
% ============================================================================
\begin{abstract}
\noindent
We show that the question whether physical reality is discrete or continuous, on its two principal readings (finite vs.\ infinite information density; computable vs.\ non-computable dynamics), cannot be decided by any finite-resource observer internal to the system. Two independent barriers are proved and machine-checked in Lean~4: a \emph{topological} barrier --- substrate discreteness is not a finitely determined (clopen) property of observable transcripts, so no finite body of evidence decides it (\leanref{Substrate.substrate\_not\_finitely\_determined}) --- and a \emph{recursion-theoretic} barrier --- no candidate law, even when supplied, can be internally certified as the law, by reduction to the halting problem (\leanref{Substrate.law\_match\_undecidable}). Under Perfect Self-Containment (PSC) these results make substrate discreteness an \emph{outsourced} property (\leanref{Substrate.substrate\_not\_internally\_decidable}), for which the only admissible adjudicator --- by a uniqueness argument this paper cites from the companion NEMS programme rather than re-derives --- is minimum-description-length (MDL) selection over complete law-sets, which selects the finite-\emph{description} class: computable transcripts together with an adjudicative extension (``glimmer'') that resolves an apparent tension with the programme's own account of physical randomness, since the purely computable class alone would contradict its determinism no-go result. We explain why Planck-scale grid searches are structurally unable to settle the question, identify asymptotic scaling exponents as the only finite observables carrying substrate information, and state the results' limits: the glimmer extension's classification is fully machine-checked; only the physical premise underlying it (that the actual Born measure constitutes complete internal semantics) and the question of which substrate actually obtains are left, honestly, open.

\NovaZenodoInAbstract
\end{abstract}

\noindent\textbf{Keywords:} substrate discreteness, computability, Baire space topology, halting problem, minimum description length, causal sets, holographic bound, Lean~4

\noindent\textbf{MSC 2020:} 03D25, 03D80, 54A05, 68Q17

\noindent\textbf{Code/Data availability:} \texttt{substrate-undecidability-lean} Lean~4 library (Mathlib only); \texttt{lake build}; repository README and module index at repository root give the full axiom audit.

% ============================================================================
% 1. INTRODUCTION
% ============================================================================
\section{Introduction}
\label{sec:intro}

``Is the universe discrete or continuous?'' sounds like one question with a definite, if
currently unknown, answer. It is not. It is at least three questions, and they are usually
run together:

\begin{itemize}[leftmargin=1.6em]
\item[\textbf{(A)}] \textbf{Information density.} Does a bounded region of space admit only
finitely many distinguishable states, or infinitely many?
\item[\textbf{(B)}] \textbf{Dynamics.} Is the time-evolution law computable (Church--Turing),
or does it access non-computable structure (hypercomputation)?
\item[\textbf{(C)}] \textbf{Chance.} Is the outcome of a measurement determined, or is there
objective randomness?
\end{itemize}

This paper's central claim is that neither reading of the question admits an internal answer:
no observer that is itself part of the physical system it is trying to characterize, and that
has only finite computational resources, can decide axis~(A) or axis~(B) --- not from evidence,
and not even if handed a candidate law outright. Both halves of that claim are proved and
machine-checked in Lean~4 against Mathlib below (\cref{sec:parti}), using nothing beyond
ordinary topology and computability theory.

A further question then arises: once no law can be internally certified, is some class of laws
nonetheless the one an internal observer should treat as evidentially favored? Answering that
second question needs one more ingredient the first does not --- \emph{Perfect
Self-Containment} (PSC), the standing closure assumption of the broader NEMS research programme
this paper draws on. PSC is not yet common vocabulary outside that programme, and the second
half of this paper's argument (Part~II, \cref{sec:mdl}) is built directly on it, so before going
any further we ground PSC carefully for a physics audience; that grounding is the rest of this
section.

\subsection{Perfect Self-Containment, for physicists}
\label{sec:psc}

Throughout, \emph{Perfect Self-Containment} (PSC) is a criterion on theories, not a hypothesis
about the cosmos. A theory satisfies PSC when its predictions are fixed by its own structure: it
needs no measure it does not supply, no vacuum selector, no external observer map, and no free
completion data to say which of its observationally inequivalent realizations is actual. The
NEMS programme proves a trichotomy --- any candidate fundamental theory is record-level
categorical, internally selecting, or not fundamental --- and a reduction showing that every
``outside dependency'' is a selection in disguise \citep{SpivackP01,SpivackP83}. Physicists
already enforce this condition without naming it: the measure problem of eternal inflation is a
problem because the measure is an external chooser \citep{Guth2007,Freivogel2011}; Wheeler's
``law without law,'' Einstein's Machian objection to absolute space (something that acts but is
not acted upon), Everett's universal wavefunction, the no-boundary proposal, and relational
quantum mechanics all exist to remove an unexplained external ingredient
\citep{Wheeler1983,Einstein1916,Everett1957,HartleHawking1983,Rovelli1996}; and the demand for no
free parameters is PSC applied to constants, with minimum-description-length selection
\citep{Rissanen1978,Solomonoff1964} as its formal version --- the same adjudication principle
this paper cites for Theorem~5 (\cref{sec:mdl}). PSC is the closure condition of which these are
instances.

Its logical status is that of an axiom, and necessarily so: a theory that derived its own
closure from a prior principle would make that principle the external selector it had just
excluded, in the same way that no consistent system certifies its own consistency
\citep{SpivackP30}. What PSC has instead is a self-consistency argument that most axioms lack.
To deny PSC is to assert an external selector \(S\). Either \(S\) is explained, in which case it
belongs to the theory and the theory is internally selecting --- PSC holds; or \(S\) is
unexplained, in which case the theory is not foundational \citep{SpivackP83}; and any further
theory that explains \(S\) faces the same dichotomy, so the regress terminates only at an
internally selecting stage \citep{SpivackP18,SpivackP23}. The denial of PSC thus posits exactly
what it declares cannot be posited. The programme's companion note on the PSC--Born fixed point
states precisely what is proved and what is assumed \citep{SpivackOnRamp}; the strongest
existence evidence is the machine-checked equivalence between PSC and the Born rule as complete
internal semantics \citep{SpivackP13,SpivackP14}, which is a fixed point rather than a
derivation from nothing. Nothing in Part~I of this paper (\cref{sec:parti}) depends on PSC;
Part~II (\cref{sec:mdl}) does, and \cref{sec:mdl} flags exactly where.

A closely related but genuinely distinct question is often folded into axis~(B) in informal
discussion: whether physical dynamics can solve computationally hard problems
\emph{efficiently} --- whether nature's time-evolution lands outside a complexity class such as
\(\mathsf{P}\) or \(\mathsf{BQP}\) using only polynomial physical resources. Call this
\textbf{axis (B\('\)): efficient computability}, the complexity class of physical dynamics.
Axis~(B\('\)) is \emph{not} axis~(B): a polynomial-time solver for an \(\mathsf{NP}\)-hard
problem is still a \emph{computable} process --- \(\mathsf{P}\) and \(\mathsf{NP}\) both sit
inside the decidable class --- so evidence bearing on axis~(B\('\)), efficiency, says nothing
about whether dynamics is computable at all, and in particular no such evidence lands anything
in \(\mathtt{Cstrict}\) (\cref{sec:model}), which requires genuine non-computability, not mere
intractability. The theorems below address axis~(B) only; axis~(B\('\)) is a
complexity-theoretic question tied to the efficient/extended Church--Turing thesis, entirely
outside their scope. We flag every place below where evidence properly belongs to
axis~(B\('\)) rather than axis~(B), since the two are conflated often enough in the literature
that keeping them apart requires active bookkeeping, not just a definition stated once.

This paper proves, and machine-checks in Lean~4 against Mathlib, that axes (A) and (B) ---
taken together as ``is the substrate discrete or continuum'' --- cannot be decided by any
observer that is itself part of the system being observed and that has only finite
resources. This holds in two independent, mutually reinforcing ways. First, no finite body of
observed data can decide it: substrate discreteness is not a \emph{finitely determined}
property of what an observer sees, in the precise sense that no prefix of observations,
however long, settles it (\cref{thm:notfindet}). Second, even supplying the observer with a
complete candidate law does not help: no internal, finite-resource procedure can certify that
the law is the correct one, on pain of deciding the halting problem (\cref{thm:lawbarrier}).
Axis (C) is orthogonal to (A) and (B); it is addressed directly by the glimmer corollary
(\cref{sec:glimmer}), whose classification is machine-checked, leaving only the Born-martingale
physical premise open (\cref{sec:open}).

The two barriers are of different character, and keeping them apart is one of the paper's
main points. The first --- call it the \emph{evidential} barrier --- is purely topological: it
is a statement about density in Baire space ($\N \to \N$ with the product topology) and does
not use computability at all; consequently no oracle, at any level of the arithmetical
hierarchy, changes a density or cardinality fact about Baire space, and nothing escapes this
barrier in the oracle sense at all. The second --- the \emph{law-based} barrier --- is
recursion-theoretic: it is a diagonal argument by reduction from the halting problem, and it
depends essentially on the observer being computable. Exactly how much extra computational
power defeats it is more subtle than ``any halting oracle'': \(\mathrm{LawMatches}(f,e)\) ---
``\(e\) matches \(f\) at every time step'' --- is a \(\Pi^0_2\) statement (a universal claim,
over every \(t\), that \(e\) halts on input \(t\) \emph{and} agrees with \(f(t)\)), and
\(\mathrm{LawMatches}\) is in fact \(\Pi^0_2\)-\emph{complete}: the reduction proving
\cref{thm:lawbarrier} gives co-halting-hardness, and a second, independent reduction
(\leanref{Substrate.law\_match\_pi2\_hard}) many-one reduces the totality problem
\(\mathrm{Tot}\) to it, which --- since \(\mathrm{Tot}\) is \(\Pi^0_2\)-complete, a standard
recursion-theoretic fact --- pins its complexity exactly at the \(\Pi^0_2\) level, strictly
above the co-halting lower bound alone. This makes the oracle-escape analysis exact rather than
one-sided: a \(\mathbf{0}''\) oracle (deciding the halting problem relative to \(\mathbf{0}'\),
sufficient for any \(\Pi^0_2\) statement) escapes \cref{thm:lawbarrier}; a plain halting
(\(\mathbf{0}'\)) oracle provably does \emph{not}, since no \(\Pi^0_2\)-complete problem is
decidable relative to \(\mathbf{0}'\). This separation (\cref{sec:whatnot}, item~5) is, as far
as we are aware, not previously stated in this form, and it matters because the two barriers
rule out the two different escape routes one might otherwise propose: ``just collect more
data'' fails against \cref{thm:notfindet}; ``but suppose the true law were simply handed to
us'' fails against \cref{thm:lawbarrier}.

Why does this matter, beyond formal tidiness? Because the standard empirical strategy for
probing substrate structure --- look for a discreteness artifact at some very small scale,
e.g.\ Lorentz-violating dispersion at the Planck length --- is not merely difficult. It is
attacking a question of the wrong shape. \cref{thm:notfindet} says that \emph{no} finite
experiment, run at any scale, with any precision, can settle discreteness in general; and the
causal-set programme (\cref{sec:history}) demonstrates concretely, for one specific
discretization scheme, that a lattice can be built with no preferred-frame artifact at all
\citep{BombelliHensonSorkin2009}. The lesson generalizes beyond that one scheme: grid searches
test whether a \emph{particular} discretization breaks a \emph{particular} symmetry, not
whether the substrate is discrete. \cref{sec:scaling} identifies what \emph{can} carry
substrate information under these barriers: not thresholds, but asymptotic scaling exponents.

Under Perfect Self-Containment (PSC) --- the criterion of \cref{sec:psc} --- the two
barriers combine into a structural diagnosis: substrate discreteness is an \emph{outsourced}
property in the sense of the Internality/Outsourcing Schema \citep{SpivackP83}. It cannot be
decided from inside by evidence or by law-certification, so if it is decided at all, the
decision procedure must live outside the space of internal tests altogether. \cref{sec:mdl}
argues, citing rather than re-deriving the uniqueness argument \citep{SpivackP34,SpivackP84},
that minimum-description-length (MDL) selection over complete law-sets is the unique
adjudicator consistent with PSC, and that it selects the finite-\emph{description} class ---
computable transcripts together with the adjudicative extension formalized in \cref{sec:glimmer}
--- which is the correct epistemic grounding for treating a compressive, finitely-described
physical theory as evidentially favored without ever claiming it has been \emph{certified}.

\paragraph{What this paper is not.} It is not a digital-physics paper. It does not argue that
the universe \emph{is} discrete, nor does it side with Zuse, Fredkin, Wolfram, or 't~Hooft
against a continuum view, nor with the opposite camp. It proves an undecidability result that
is agnostic about which substrate obtains, and it explains why every position in that debate is
structurally unprovable from the inside --- a claim that is stronger than, and different in
kind from, taking a side (\cref{sec:discussion}).

\paragraph{Structure.} \cref{sec:psc} grounds Perfect Self-Containment for a physics audience,
before it is invoked later in this section and in Part~II. \cref{sec:history} traces the
intellectual history of this question, from the classical continuum assumption through
computability theory to the modern empirical programmes, and situates this paper's result
relative to all of them. \cref{sec:axes} states the three axes and the existing physical
evidence bearing on each. \cref{sec:model} gives the
formal model (worlds, transcripts, internal tests, finite determination), matching
\leanref{Substrate/Basic.lean}. \cref{sec:parti} states and sketches the proofs of the Part~I
theorems. \cref{sec:whatnot} states carefully what the theorems do \emph{not} say.
\cref{sec:scaling} locates the only finite observables that carry substrate information.
\cref{sec:mdl} states Theorem~5 (outsourcing and MDL) and, in \cref{sec:glimmer}, the glimmer
corollary that settles axis~C's classification. \cref{sec:discussion} is a synthesizing
discussion. \cref{sec:open} lists the remaining open problems, including the Born-martingale
lemma for axis~C. \cref{sec:conclusion} concludes. Appendix~\ref{app:lean} lists the Lean
declarations; Appendix~\ref{app:status} gives the epistemic status table.

% ============================================================================
% 2. HISTORICAL CONTEXT
% ============================================================================
\section{Historical Context: A Question Posed and Re-Posed}
\label{sec:history}

The discrete/continuum question has been asked in substantially different vocabularies by
each generation that has taken it up, and the vocabulary shift is not cosmetic --- each new
formal tool changed what counted as an answer, and, we will argue, each generation's tools
eventually ran into some version of the same wall this paper formalizes.

\subsection{The continuum default, and its first crack}

Classical mechanics inherited its substrate assumption from geometry rather than from physics:
space and time were modeled on the real continuum because that was the available mathematics,
not because the continuum had been empirically established. Newtonian dynamics, and later
field theory through general relativity, is written on differentiable manifolds with no
built-in minimum length or minimum state count. For roughly three centuries, ``the universe is
continuous'' was less a hypothesis than an unexamined modeling default, and the question this
paper asks --- can an internal observer decide it? --- had no formal handle, because there was
no rigorous notion yet of what a finite-resource decision procedure even was.

That handle appeared with computability theory in the 1930s (Turing, Church, Kleene), but its
relevance to physics was not immediate. The connection was made explicit by Robin Gandy, who
in \citeyear{Gandy1980} asked, in effect, what physical principles a machine's operation must
satisfy for the Church--Turing thesis to extend from ``mechanisms'' in the abstract to devices
built from actual matter \citep{Gandy1980}. Gandy's analysis --- bounding the local density of
information and the propagation speed of causal influence --- was the first serious attempt to
tie ``how much information can be packed into a bounded physical region'' (our axis~A) to a
computability claim about physical process (our axis~B). It is a direct ancestor of both
Bekenstein's later entropy bound and of the present paper's insistence that axes~A and~B, while
conceptually distinct, are both instances of one underlying phenomenon: an internal observer's
access to the substrate is mediated entirely by a finite-alphabet readout channel, and every
question about the substrate is really a question about what can be inferred from that channel.

\subsection{The digital-physics tradition}

Independently of Gandy, and starting earlier, a distinct tradition proposed that physical
reality is not merely finite-information but literally computational in its substrate. Konrad
Zuse's 1969 \emph{Rechnender Raum} (``Calculating Space'') \citep{Zuse1969,Zuse1970} is
usually credited as the founding text: Zuse conjectured that the universe is being computed by
a discrete, cellular-automaton-like process, decades before ``digital physics'' had a name and
before Bekenstein or Bousso had given any information-theoretic content to ``finite substrate
information.'' Edward Fredkin developed this into a research programme in the 1980s,
proposing that reversible universal cellular automata are the correct microphysical model
\citep{Fredkin1990}, and Stephen Wolfram's \emph{A New Kind of Science}
\citep{Wolfram2002} pushed the same idea into a large-scale programme of searching cellular
automaton rule space for candidates that reproduce known physics. Gerard~'t~Hooft took a
related but distinct route from a different motivation: dissatisfied with the interpretive
status of quantum mechanics, he proposed that quantum indeterminism is an emergent, coarse-grained
description of an underlying deterministic, discrete (cellular-automaton-like) dynamics
\citep{tHooft1999,tHooft2016} --- so 't~Hooft's discreteness claim is entangled with a
specific stance on axis~C (determinism) that Zuse, Fredkin, and Wolfram's proposals do not
require.

These positions share a structure worth naming precisely, because it is exactly the structure
this paper shows cannot be internally vindicated: each proposes a \emph{specific} discrete law
or class of discrete laws and treats successfully reproducing observed physics as evidence, or
even proof, that the substrate literally is of that computational character. \cref{thm:lawbarrier}
below says that no internal, finite-resource procedure can certify that a candidate law is the
one generating the observed transcript forever, however well it has matched so far. This is not
a criticism specific to the digital-physics tradition --- the identical epistemic ceiling applies
to any continuum theory claiming certification, including general relativity taken at face
value as a claim about the actual substrate. It is a symmetric barrier, and one purpose of this
paper's abstraction (worlds, transcripts, tests, defined independent of any particular physical
theory) is to make that symmetry a theorem rather than an assertion.

\subsection{Complications from the computability side}

The digital-physics tradition's implicit dichotomy --- discrete-and-computable versus
continuum-and-mysterious --- did not survive contact with two independent developments inside
computability theory and analysis themselves, both of which this paper's model deliberately
inherits rather than glosses over.

First, David Deutsch's \citeyear{Deutsch1985} formulation of the \emph{physical} Church--Turing
principle --- that every finitely realizable physical system can be perfectly simulated by a
universal computing machine operating by finite means \citep{Deutsch1985} --- was explicitly a
physical claim, not a mathematical theorem, and Deutsch's own paper already flagged the
subtlety that a computable law can still support qualitatively new physical concepts (a
quantum computer efficiently doing things intractable classically) without leaving the
computable realm. The physical Church--Turing thesis being an empirically contentful claim
about physical law, rather than a mathematical certainty, is precisely the kind of thing
\cref{thm:lawbarrier} shows cannot be internally certified.

Second, and more corrosive to the naive dichotomy, Hava Siegelmann's analysis of analog
recurrent neural networks with real-valued (non-rational) weights showed that such networks
can, in principle, compute functions no Turing machine can compute, evaluated over unbounded
time \citep{SiegelmannSontag1994,Siegelmann1999}. This is a genuine, mathematically rigorous
hypercomputation candidate that lives entirely within continuum mathematics and requires no
appeal to exotic new physics --- it only requires that some physical quantity actually realize
an uncomputable real number to unbounded precision. B.~Jack Copeland's philosophical taxonomy
of hypercomputation \citep{Copeland2002,Copeland2002b} carefully separates several logically
distinct routes to physical hypercomputation (accelerating Turing machines, infinite-precision
analog computation, quantum hypercomputation proposals, oracle-consulting machines) that are
frequently conflated in informal discussion; that taxonomy is one of the reasons this paper
insists on stating precisely, and separately, what each of its two barriers does and does not
assume about the observer's computational resources (\cref{sec:whatnot}, item~5).

Third, and most directly relevant to this paper's modeling choice for the continuum class,
Marian Pour-El and Ian Richards proved that the classical wave equation, with perfectly
computable initial data, can have a unique solution that is \emph{not} computable
\citep{PourElRichards1981,PourElRichards1989}. (A regularity caveat Penrose raised
historically, and any careful referee will too: the non-computable solution arises from
initial data that is computable but not \(C^1\); with smooth, \(C^1\)-computable data the
solution \emph{is} computable. This does not affect how the result is used below ---
\cref{rem:modeling} and \cref{sec:open} only need the existence of \emph{some} non-computable
member of the computable-analysis-refined continuum class, which the non-\(C^1\) construction
still supplies.) This result is the single clearest formal
demonstration that ``continuum substrate'' does not entail ``hypercomputational substrate''
and, symmetrically, that ``discrete substrate'' is not needed to obtain ``computable
substrate'' --- computability and cardinality of the state space are logically independent
axes. It undercuts any easy identification of continuum with dangerous-and-uncomputable or
discrete with safe-and-tame, and it is exactly why this paper's formal model
(\cref{sec:model}) defines the discrete/continuum distinction by the cardinality and structure
of the state space (\leanref{Substrate.DiscreteWorld}, \leanref{Substrate.ContinuumWorld}),
not by whether the resulting transcript is computable, and why \(\mathtt{Cstrict}\) --- the
class of continuum transcripts that happen to be non-computable --- is defined as a strict
subset of the continuum class rather than identified with it. A computable-analysis refinement
of the continuum class (restricting to Cauchy-modulus-computable reals) would shrink
\(\mathtt{Cstrict}\); by the Pour--El--Richards result it would not empty it, because their
non-computable solution already arises from computable initial data under exactly that kind of
restriction. We return to this in \cref{sec:whatnot}, item~4.

\subsection{The empirical side: how physics actually tried to test axis A}

While the digital-physics tradition and the computability-theoretic complications developed
largely in parallel, mainstream theoretical physics developed its own, independent answer to
axis~A --- not by conjecturing a specific discrete mechanism, but by deriving an information
bound from established physics. Jacob Bekenstein's \citeyear{Bekenstein1981} universal bound
--- that the entropy-to-energy ratio of any bounded physical system is bounded above by a
quantity proportional to its radius, \(S \le 2\pi k_B E R / \hbar c\) --- \citep{Bekenstein1981}
was derived from black-hole thermodynamics and the second law, not from any discreteness
conjecture, and already gives a positive answer to axis~A for systems of bounded energy: only
finitely many distinguishable states fit in a bounded region of given energy content. Gerard
't~Hooft \citeyear{tHooft1993} and Leonard Susskind \citeyear{Susskind1995} then proposed the
qualitatively stronger \emph{area} bound --- that the maximum entropy (equivalently,
information content) of any bounded region is proportional to its \emph{boundary area}, not
its volume, at roughly one bit per Planck area --- building explicitly on Bekenstein's result
but going beyond it: the area bound is a claim about geometry itself, not merely a bound tied
to a system's energy content \citep{tHooft1993,Susskind1995}. Raphael Bousso generalized the
area bound into a fully covariant form applicable to arbitrary spacetimes and causal structures
\citep{Bousso1999,Bousso2002}. This is the historically important case of axis~A being
addressed \emph{without} committing to any particular discretization scheme, computational
interpretation, or even to the claim that spacetime itself is discrete --- the holographic
bound constrains information content, not geometry.

The natural next question --- if information density is finite, does spacetime itself have to
be a discrete lattice, and if so, wouldn't such a lattice show up experimentally as a preferred
frame (Lorentz violation)? --- is where the causal-set programme becomes historically
essential. Building on ideas that trace back through Bombelli, Lee, Meyer, and Sorkin's 1987
proposal that spacetime is fundamentally a discrete partial order (a causal set)
\citep{Bombelli1987}, Bombelli, Henson, and Sorkin proved in \citeyear{BombelliHensonSorkin2009}
that a discrete causal set, generated by a Poisson sprinkling process, can be constructed to be
exactly Lorentz invariant \emph{in distribution} --- no preferred rest frame, no direction
picked out by the discretization \citep{BombelliHensonSorkin2009}. This is a narrower and more
precise claim than ``no discreteness artifact whatsoever'': causal sets built this way still
exhibit their own discreteness signatures (non-locality effects, so-called ``swerves'' in
particle trajectories), which are themselves Lorentz-invariant rather than absent. The theorem
rules out a preferred \emph{frame}, not discreteness signatures altogether --- a distinction
worth stating precisely, because it makes this paper's point sharper rather than weaker: it
shows discreteness need not manifest as a symmetry violation at all, which is exactly why a
naive ``look for the grid'' experiment is attacking a question of the wrong shape even before
\cref{thm:notfindet}'s general obstruction is invoked. This result mattered historically
because it dissolved what had looked like a knock-down objection to any discrete spacetime
theory (``a lattice necessarily picks out a preferred frame, we would have seen it by now''),
and thereby made discrete-spacetime theories respectable candidates again.

We engage this history directly because it is the single clearest historical instance of
exactly the phenomenon this paper's Theorem~2 (\cref{thm:notfindet}) proves in general. The
naive Lorentz-violation test is precisely a finitely-determined property attempt: look for a
specific, bounded-resource observable signature (anisotropic dispersion, a preferred direction)
that a specific discretization scheme would produce. The causal-set programme's own resolution
--- construct the discretization to be exactly Lorentz-invariant --- is a demonstration, for
one concrete scheme, that \emph{that particular} finite test cannot distinguish the discrete
case from the continuum case. \cref{thm:notfindet} generalizes this from ``this one clever
scheme evades this one clever test'' to ``no finite-resource test whatsoever, against any
discretization scheme whatsoever, can settle the question, because the two transcript classes
are dense into each other in Baire space.'' The causal-set community's Lorentz-invariance
achievement is therefore best read, in retrospect, as a special-case patch that happens to
defeat one attack; it was never, and could never have been on its own terms, a general
resolution of the underlying undecidability, because no finite construction can be a general
resolution of a statement about the density of two disjoint sets in an infinite product space.

\subsection{Modern empirical programmes on axis B\texorpdfstring{\('\)}{'}, not axis B}

Parallel to the causal-set programme's engagement with axis~A, two modern lines of work have
tried to test whether physical dynamics can defeat known complexity-class bounds --- axis~B\('\)
as introduced in \cref{sec:intro}, not axis~B itself, since neither line of work bears on
whether dynamics is \emph{computable} at all, only on how \emph{efficiently} computable it is.
Scott Aaronson's ``NP-complete problems and physical reality'' \citep{Aaronson2005} surveys and
evaluates every serious proposal (from relativistic time dilation near a black hole to exotic
quantum-gravitational effects) for physically solving NP-hard problems in polynomial time,
concluding that none survives scrutiny and proposing, in effect, that nature's apparent refusal
to solve NP-hard problems efficiently is itself a physical principle worth taking as seriously
as the second law of thermodynamics. Seth Lloyd's estimates of the computational capacity of
the observable universe \citep{Lloyd2002,Lloyd2006} --- bounding the number of elementary
logical operations and bits the universe could have performed over its history, using the same
Bekenstein-type bounds --- are the natural quantitative complement: they do not test
hypercomputation (axis~B) at all, but put a number on how much classical computation ``fits,''
which is the natural backdrop against which any claimed axis-B\('\) anomaly (a physical process
solving something outside \(\mathsf{BQP}\) with polynomial resources) would have to be
evaluated. Both of these are external-literature results about axis~B\('\); the programme's own
internal result on axis~B \emph{itself} --- not B\('\) --- is NEMS Paper~35, ``Oracles as
External Selectors,'' which proves a hypercomputation-exclusion taxonomy (no admissible
internal procedure consults an oracle) directly from the Internality/Outsourcing Schema
\citep{SpivackP35}; \cref{thm:certifier} below complements it with a Lean-checked,
model-theoretic certifier-silence result reached independently and by a different route.

Both Aaronson's and Lloyd's results are, in our terminology, attempts to test axis~B\('\) empirically by looking for a
single dramatic counterexample or bound violation; neither bears on axis~B as this paper defines
it, since a fast NP-hard solver is still a computable process. \cref{thm:certifier} below
(Theorem~3, sound certifier never fires) is nonetheless the right formal lens for how such a
result should be interpreted even if it were found, on its own complexity-theoretic terms: by
the same prefix-density argument (\cref{thm:prefix}) that proves \cref{thm:certifier} for the
discrete/continuum case, any finite observed prefix of ``evidence for the violation'' is equally
consistent with an ordinary, bound-respecting computable continuation, so no internal test built
from that finite prefix could be sound and ever fire on this kind of claim either --- the formal
reason a single demonstration is never internally certifiable, not the informal worry that it
might be a resource-hiding artifact. Only the \emph{scaling} of physical resource against
instance size, across a family of instances, can be reported as evidence --- never certified as
proof. This is why \cref{sec:scaling} treats scaling exponents, not individual demonstrations,
as the correct unit of empirical content for axis~B\('\), continuing the same methodological
lesson the causal-set programme taught for axis~A.

\subsection{Where this paper sits}

Every result surveyed above is either (i) a specific mechanism proposal (digital physics), (ii)
a specific complication showing the discrete/continuum dichotomy does not cleanly track
computability (Deutsch, Siegelmann, Pour-El--Richards), (iii) a specific information bound
derived from established physics (Bekenstein, Bousso), or (iv) a specific empirical test or
estimate (causal sets' Lorentz invariance, Aaronson's survey, Lloyd's estimates). This paper is
none of these. It does not propose a mechanism, does not derive a new physical bound, and does
not run or evaluate an empirical test. It proves, for an observer defined abstractly enough to
cover every mechanism proposal and every empirical test above as a special case, that the
entire \emph{class} of such tests --- not any one of them, all of them, uniformly --- cannot
settle the question. That is a claim about the logical structure of the problem, not a new
entrant into the substantive debate about which substrate is correct, and it is the reason the
paper is agnostic between Zuse/Fredkin/Wolfram/'t~Hooft on one side and an unrestricted
continuum view on the other: both sides are making a claim this paper proves can never be
internally certified, regardless of which of them happens to be true.

% ============================================================================
% 3. THE THREE AXES AND PRIOR EVIDENCE
% ============================================================================
\section{The Three Axes and Prior Evidence}
\label{sec:axes}

\begin{table}[htbp]
\centering
\small
\begin{tabular}{@{}>{\raggedright\arraybackslash}p{1.7cm}>{\raggedright\arraybackslash}p{2.8cm}>{\raggedright\arraybackslash}p{2.8cm}>{\raggedright\arraybackslash}p{3.9cm}>{\raggedright\arraybackslash}p{3.5cm}@{}}
\toprule
\textbf{Axis} & \textbf{Discrete pole} & \textbf{Continuum pole} & \textbf{Existing evidence} & \textbf{Attack surface} \\
\midrule
A. Information density & Finite states per bounded region & Infinitely many distinguishable states & Bekenstein's energy bound + 't~Hooft\slash Susskind\slash Bousso's holographic area bound \(\Rightarrow\) finite \citep{Bekenstein1981,tHooft1993,Susskind1995,Bousso1999,Bousso2002} & Exactness of the bound; metrological scaling \\
B. Dynamics & Computable (Church--Turing) & Non-computable / hypercomputation accessible & Siegelmann's analog RNN construction shows hypercomputation is mathematically consistent \citep{SiegelmannSontag1994,Siegelmann1999}; no confirmed physical instance & \cref{thm:lawbarrier}: no internal law-certification, whether the candidate law is computable or not \\
B\('\). Efficient computability \emph{(subsidiary to B --- complexity, not computability; see \cref{sec:intro})} & \(\mathsf{P}\)/\(\mathsf{BQP}\)-computable with polynomial resources & \(\mathsf{NP}\)-hard (or harder) solved with polynomial resources & Aaronson: nature refuses \(\mathsf{NP}\)-hard in poly time \citep{Aaronson2005}; Lloyd's computational-capacity bounds \citep{Lloyd2002,Lloyd2006} & Complexity class of physical computation --- outside the Lean theorems below, which address computability, not efficiency \\
C. Chance & Deterministic & Objective randomness & Bell's theorem, assuming no superdeterminism \citep{Bell1964}; Hossenfelder--Palmer's case for taking superdeterminism seriously is the standing challenge to that assumption, not support for it \citep{HossenfelderPalmer2020} & Interacts with B via the ``glimmer'' programme (\(\Omega\)-like traces at degree \(\mathbf{0}'\)); \cref{sec:open} \\
\bottomrule
\end{tabular}
\caption{The axes conflated by ``is the universe discrete or continuous.'' This paper's Part~I theorems (\cref{sec:parti}) show axes A and B, \emph{taken together} as the discrete-versus-strict-continuum distinction (\cref{sec:model}), are internally undecidable; \cref{rem:axisA} gives a separate, axis-A-specific dense pair (\(\mathtt{Dfin}\), eventually periodic sequences) to which the same mechanism applies directly, machine-checked (\cref{rem:axisA}). Axis B\('\) is a distinct, complexity-theoretic question (efficiency, not computability) that the theorems do not address and that is listed here only to correctly separate it from axis B's existing evidence. Axis C is orthogonal and is addressed by the glimmer corollary (\cref{sec:glimmer}): its definitions and separations are machine-checked, with one bridge lemma (the Born-martingale completeness premise) disclosed as open.}
\label{tab:axes}
\end{table}

Axes~A and~B are exactly the two the theorems below show to be internally undecidable; axis~C
is genuinely different in kind and is addressed only by the glimmer corollary
(\cref{sec:glimmer}), whose definitions and separations are machine-checked but whose one
remaining bridge lemma is disclosed as open, not folded in as an established result. Axis~B\('\)
(\cref{tab:axes}) is also not addressed: it is a complexity-theoretic question about efficiency,
not a computability question, and lies outside what \cref{thm:lawbarrier} or any other theorem
below speaks to. The Bombelli--Henson--Sorkin causal-set result
\citep{BombelliHensonSorkin2009} is the standing witness that discreteness on axis~A leaves no
Lorentz-violating grid artifact, so the traditional Planck-scale experimental programme tests a
particular scheme's symmetry properties, not substrate as such (\cref{sec:history}).

\begin{remark}[Axis A is not literally \(\D\), and has its own dense pair]
\label{rem:axisA}
\(\D\) as defined (\cref{def:classes}) is an axis-B class --- computability of transcript --- not
literally axis~A's ``finite states per bounded region.'' A finite-state world with a
non-computable transition rule has finite information density (axis A) but a non-computable
transcript, and a continuum world with a computable rule has the reverse; \(\D\) tracks neither
in isolation. To make the abstract's claim to cover ``both principal readings'' fully formal
rather than only evidentially motivated (\cref{tab:axes}, row A), axis~A gets its own dense pair
and its own instance of \cref{thm:notfindet}: let \(\mathtt{Dfin} := \{f \mid \exists\, w
\text{ finite-state},\ w.\mathrm{transcriptFn} = f\}\) be the transcripts of worlds with finite
state type. \(\mathtt{Dfin}\) coincides exactly with the eventually periodic sequences
(\leanref{Substrate.Dfin\_iff\_eventuallyPeriodic}) --- a pigeonhole argument on the finitely many
reachable states forces eventual periodicity, and conversely every eventually periodic sequence
is realized by a finite-state world built from its pre-period and period. \(\mathtt{Dfin}\) is
dense (\leanref{Substrate.Dfin\_dense}), its complement is dense (\leanref{Substrate.Dfin\_compl\_dense}),
and the two are disjoint by construction, so the abstract dense-pair form of Theorem~2
(\leanref{Substrate.not\_finitely\_determined\_of\_dense\_pair}, module \leanref{Substrate/NotFinDet.lean})
applies verbatim to this different pair:
\[
  \neg\, \exists\, P,\ \mathrm{FinitelyDetermined}(P) \,\wedge\, (\forall f \in \mathtt{Dfin},\
  P(f)) \,\wedge\, (\forall f \in \mathtt{Dfin}^c,\ \neg P(f))
\]
(\leanref{Substrate.axisA\_not\_finitely\_determined}, module \leanref{Substrate/AxisA.lean}).
Axis~A is therefore internally undecidable in its own right, by the identical topological
mechanism as axes A and B taken together, and not merely by inheritance from axis~B's \(\D\).

\smallskip\noindent\emph{Lean status:} \leanref{Substrate/AxisA.lean} is machine-checked
--- status \A. It is part of this repository's green \leanref{Substrate.lean} import (\texttt{lake
build}: 8262/8262 jobs, zero errors, zero \texttt{sorry}), and every theorem named above is
confirmed by direct \texttt{\#print axioms} to depend on at most the three standard Mathlib
axioms (\texttt{propext}, \texttt{Classical.choice}, \texttt{Quot.sound}); see
Appendix~\ref{app:lean} for the full declaration table.
\end{remark}

\FloatBarrier

% ============================================================================
% 4. FORMAL MODEL
% ============================================================================
\section{Formal Model}
\label{sec:model}

The formal development (Lean module \leanref{Substrate/Basic.lean}) fixes an abstraction of
``physical system with an internal observer'' deliberately thin enough to cover every substrate
proposal in \cref{sec:history} as a special case, and precise enough to support theorems.

\subsection{Worlds and transcripts}

\begin{definition}[World]
\label{def:world}
A \emph{world} (\leanref{Substrate.World}) is a record consisting of a state type \(S\), an
initial state \(\mathrm{init} : S\), a deterministic step function \(\mathrm{step} : S \to S\),
and a finite-alphabet observer readout \(\mathrm{read} : S \to \N\).
\end{definition}

The readout is where ``internal observer'' enters formally, and it is the single most important
modeling choice in the paper. Whatever the substrate actually is --- finite state machine,
continuum field configuration, or anything else --- the observer receives only a stream of
natural numbers, one symbol per time step. Everything an internal observer can possibly know
about the substrate factors through this stream.

\begin{definition}[Transcript]
\label{def:transcript}
The \emph{transcript} of a world \(w\) (\leanref{Substrate.World.transcript},
\leanref{Substrate.World.transcriptFn}) is the function \(\N \to \N\) given by
\[
  w.\mathrm{transcript}(t) \;=\; w.\mathrm{read}\bigl(w.\mathrm{step}^{[t]}(w.\mathrm{init})\bigr).
\]
\end{definition}

The transcript is the \emph{only} thing the observer ever sees; every downstream statement in
this paper is a statement about elements of Baire space \(\N \to \N\).

\subsection{Two substrate classes, defined by what they can produce}

Substrate classes are defined by the transcripts they can produce, not by the state type
directly --- this is what makes the indistinguishability theorems (\cref{sec:parti}) clean, and
it is a deliberate anti-smuggling choice: a definition of ``discrete'' phrased in terms of the
state type alone would risk building the conclusion into the premise.

\begin{definition}[Discrete and continuum worlds]
\label{def:classes}
A world \(w\) is \emph{discrete} (\leanref{Substrate.DiscreteWorld}) if its state type admits a
\texttt{Primcodable} structure and both \(\mathrm{step}\) and \(\mathrm{read}\) are
\texttt{Computable}. A world \(w\) is a \emph{continuum} world
(\leanref{Substrate.ContinuumWorld}) if its state type is isomorphic to \(\mathrm{Fin}\,n \to
\R\) for some \(n\), with \emph{arbitrary} (unrestricted classical) dynamics.
\[
  \D := \{f \mid \exists\, w \text{ discrete},\ w.\mathrm{transcriptFn} = f\}
  \qquad
  \mathtt{C} := \{f \mid \exists\, w \text{ continuum},\ w.\mathrm{transcriptFn} = f\}
\]
\[
  \mathtt{Cstrict} := \mathtt{C} \setminus \{f \mid \mathrm{Computable}\ f\}
\]
(\leanref{Substrate.D}, \leanref{Substrate.C}, \leanref{Substrate.Cstrict}.)
\end{definition}

\begin{remark}[Modeling honesty]
\label{rem:modeling}
\(\mathtt{ContinuumWorld}\) places no computability restriction whatsoever on \(\mathrm{step}\);
consequently \(\mathtt{C}\) turns out to be \emph{all} of Baire space
(\leanref{Substrate.mem\_C}, \cref{sec:parti}): an unrestricted continuum with arbitrary
classical dynamics can realize any readout stream. This is a deliberate and disclosed modeling
choice, not an oversight --- it reflects the historical fact (\cref{sec:history}) that
``continuum'' as classically conceived carries no built-in computability restriction, and it is
exactly why the interesting content of the continuum class lives in \(\mathtt{Cstrict}\), the
sub-class that is \emph{also} non-computable. A computable-analysis refinement of
\(\mathtt{ContinuumWorld}\) --- restricting \(\mathrm{step}\) to act by Cauchy-modulus-computable
maps on Cauchy-modulus-computable reals --- would shrink \(\mathtt{Cstrict}\), but by
Pour-El--Richards \citep{PourElRichards1981} would not empty it, since their non-computable
wave-equation solution already arises from computable initial data under exactly that kind of
restriction. Mathlib does not currently provide computable analysis, so this refinement is left
for future work (\cref{sec:open}); the theorems below do not depend on it.
\end{remark}

\subsection{Internal tests and finite determination}

\begin{definition}[Prefix, internal test]
\label{def:test}
For \(f : \N \to \N\) and \(n \in \N\), \(\mathrm{Prefix}(f,n)\) (\leanref{Substrate.Prefix}) is
the list of the first \(n\) values of \(f\). An \emph{internal test}
(\leanref{Substrate.InternalTest}) is a pair of a budget \(n \in \N\) and a computable verdict
function \(\mathrm{verdict} : \mathrm{List}\,\N \to \mathrm{Bool}\); it runs
(\leanref{Substrate.InternalTest.run}) as \(T.\mathrm{run}(f) := T.\mathrm{verdict}(\mathrm{Prefix}(f,
T.\mathrm{budget}))\).
\end{definition}

\begin{definition}[Finitely determined]
\label{def:findet}
A property \(P\) of transcripts is \emph{finitely determined}
(\leanref{Substrate.FinitelyDetermined}) if some prefix length \(n\) settles it for every pair
of transcripts agreeing on that prefix:
\[
  \exists\, n,\ \forall f\,g,\ \mathrm{Prefix}(f,n) = \mathrm{Prefix}(g,n) \Rightarrow (P(f) \leftrightarrow P(g)).
\]
Equivalently, \(P\) is clopen in the product topology on \(\N \to \N\).
\end{definition}

An internal test decides a property exactly when the property is internally decidable
(\cref{sec:mdl}, \leanref{Substrate.InternallyDecidable}); every internally decidable property
is finitely determined (\leanref{Substrate.internally\_decidable\_implies\_finitely\_determined}),
since a test only ever inspects a bounded-length prefix.

% ============================================================================
% 5. THEOREMS, PART I
% ============================================================================
\section{Theorems, Part I: The Two Barriers}
\label{sec:parti}

All results in this section are proved and machine-checked (status \A) in the Lean module
family \leanref{Substrate/}, against Mathlib only.

\begin{lemma}[Transcript characterisation of \(\D\)]
\label{lem:transcript-char}
\(f \in \D \iff \mathrm{Computable}\ f\).
{\raggedright\emph{(\leanref{Substrate.transcript\_characterisation}, module \leanref{Substrate/Discrete.lean}.)}\par}
\end{lemma}

\begin{proof}[Proof sketch]
(\(\Rightarrow\)) If \(w\) is discrete, iterating a computable \(\mathrm{step}\) from a
computable state is computable (\leanref{Substrate.World.iterate\_computable}, via
\texttt{Computable.nat\_rec}), and composing with computable \(\mathrm{read}\) shows the
transcript is computable (\leanref{Substrate.World.transcript\_computable\_of\_discrete}).
(\(\Leftarrow\)) Given computable \(f\), take \(S := \N\), \(\mathrm{step} :=
\mathrm{succ}\), \(\mathrm{init} := 0\), \(\mathrm{read} := f\); this world is discrete by
construction and its transcript is \(f\) by \(\mathrm{succ}^{[t]}(0) = t\).
\end{proof}

\begin{lemma}[Embedding]
\label{lem:embedding}
Every discrete world's transcript is realized by \emph{some} continuum world; hence \(\D
\subseteq \mathtt{C}\). \emph{(\leanref{Substrate.discrete\_subset\_continuum}, module
\leanref{Substrate/Embed.lean}.)}
\end{lemma}

\begin{proof}[Proof sketch]
This is the immediate specialisation of the stronger fact, proved directly, that \emph{every}
\(f : \N \to \N\) whatsoever lies in \(\mathtt{C}\) (\leanref{Substrate.mem\_C}): since \(\N \to
\N\) and \(\R\) have the same cardinality (\texttt{Cardinal.mk} agreement, via
\(\aleph_0^{\aleph_0} = \mathfrak{c}\)), \texttt{Mathlib}'s \texttt{Cardinal.eq} converts that
cardinality equality into a \emph{classical}, non-constructive equivalence \((\N \to \N) \simeq
(\mathrm{Fin}\,1 \to \R)\) --- this step is where \texttt{Classical.choice} (and, via
\texttt{Cardinal.mk}'s quotient construction, \texttt{Quot.sound}) does genuine work in this
proof, not merely disclosed bookkeeping: no \emph{explicit}, computable equivalence between the
two types is exhibited or exhibitable by this argument. Given that equivalence, take the
continuum world's state type to be \(\N \to \N\) itself, \(\mathrm{init} := f\),
\(\mathrm{step}\) the left-shift, and \(\mathrm{read}\) the head symbol; its transcript is
exactly \(f\).
\end{proof}

\begin{theorem}[Prefix Indistinguishability]
\label{thm:prefix}
For every finite list \(p\), there exist \(d \in \D\) and \(c \in \mathtt{Cstrict}\), both
extending \(p\) as a prefix: \(d\) is \(p\) followed by zeros (computable), and \(c\) is
\(p\) followed by a fixed non-computable tail built from the halting problem. \emph{(\leanref{Substrate.prefix\_indistinguishable},
module \leanref{Substrate/Prefix.lean}.)}
\end{theorem}

\begin{proof}[Proof sketch]
The discrete witness is the zero-extension of \(p\), computable by construction
(\leanref{Substrate.prefix\_getD\_eq}). The strict-continuum witness agrees with \(p\) on the
prefix and then continues with a fixed function \(\mathrm{haltingIndicator} : \N \to \N\)
(\leanref{Substrate.haltingIndicator}), defined as the (classically well-defined, non-effective)
indicator of whether a code halts on input \(0\); this function is provably not computable
(\leanref{Substrate.not\_computable\_haltingIndicator}), by reduction from
\texttt{ComputablePred.halting\_problem}, and this non-computability transfers to the
prefix-extended witness. By \cref{lem:embedding} the witness also lies in \(\mathtt{C}\), hence
in \(\mathtt{Cstrict}\).
\end{proof}

\begin{corollary}
\label{cor:disjoint-dense}
\(\D\) and \(\mathtt{Cstrict}\) are disjoint (\leanref{Substrate.D\_disjoint\_Cstrict}, an
immediate consequence of \cref{lem:transcript-char}) and both dense in Baire space (immediate
from \cref{thm:prefix}, since every finite prefix extends into both classes).
\end{corollary}

\begin{theorem}[Substrate is not finitely determined]
\label{thm:notfindet}
No finitely determined property of transcripts holds on all of \(\D\) and fails on all of
\(\mathtt{Cstrict}\):
\[
  \neg\, \exists\, P,\ \mathrm{FinitelyDetermined}(P) \,\wedge\, (\forall f \in \D,\ P(f))
  \,\wedge\, (\forall f \in \mathtt{Cstrict},\ \neg P(f)).
\]
{\raggedright\emph{(\leanref{Substrate.substrate\_not\_finitely\_determined}, module \leanref{Substrate/NotFinDet.lean}.)}\par}
\end{theorem}

\begin{proof}[Proof sketch]
Immediate from \cref{thm:prefix}: if \(P\) were settled by prefix length \(n\), apply
\cref{thm:prefix} to the all-zero prefix of length \(n\) to get \(d \in \D\) and \(c \in
\mathtt{Cstrict}\) agreeing on that prefix; then \(P(d) \leftrightarrow P(c)\) contradicts \(P(d)\)
true and \(P(c)\) false.
\end{proof}

\begin{keybox}[title={This is the headline result --- and it does not need the halting problem}]
Theorem~2 is pure topological density in Baire space. Its proof invokes \cref{thm:prefix},
whose \(\mathtt{Cstrict}\) witness happens to be built from a halting-problem construction, but
nothing about Theorem~2 \emph{itself} depends on recursion theory: any two dense, disjoint
subsets of \(\N \to \N\) would give the same conclusion. This is a Lean-checked fact, not
only a proof-sketch observation: the abstract form
\[
  (\forall p,\ \exists a \in A,\ \mathrm{Prefix}(a,|p|) = p) \to (\forall p,\ \exists b \in B,\
  \mathrm{Prefix}(b,|p|) = p) \to \mathrm{Disjoint}\ A\ B \to {}
\]
\[
  \neg\, \exists\, P,\ \mathrm{FinitelyDetermined}(P) \wedge (\forall f \in A,\ P(f)) \wedge
  (\forall f \in B,\ \neg P(f))
\]
is proved directly from prefix combinatorics, with no reference to \texttt{Computable} or
\texttt{Nat.Partrec} anywhere in its proof term
(\leanref{Substrate.not\_finitely\_determined\_of\_dense\_pair}, module
\leanref{Substrate/NotFinDet.lean}; \texttt{\#print axioms} shows only the standard
\texttt{propext, Classical.choice, Quot.sound}), and \cref{thm:notfindet} itself is derived from
it by instantiating \(A := \D\), \(B := \mathtt{Cstrict}\) via \cref{thm:prefix}
(\leanref{Substrate.substrate\_not\_finitely\_determined}). The barrier is topological
\emph{before} it is recursion-theoretic, and Lean shows this at the level of the proof term,
not only the theorem statement. This separation is developed further in \cref{sec:whatnot},
item~5, and is one of the two independent barriers advertised in the abstract.
\end{keybox}

\begin{theorem}[No sound internal certifier fires]
\label{thm:certifier}
Let \(T\) be an internal test.
\begin{enumerate}[label=(\alph*), leftmargin=1.6em]
\item If \(T\) is \emph{sound for continuum}
(\leanref{Substrate.SoundForContinuum}: \(T.\mathrm{run}(f) = \mathrm{true} \Rightarrow f \in
\mathtt{Cstrict}\)), then \(T\) never returns \(\mathrm{true}\)
(\leanref{Substrate.sound\_for\_continuum\_never\_fires}).
\item Symmetrically, if \(T\) is
\emph{sound for discreteness} (\leanref{Substrate.SoundForDiscreteness}), \(T\) never returns
\(\mathrm{true}\) (\leanref{Substrate.sound\_for\_discreteness\_never\_fires}).
\end{enumerate}
{\raggedright\emph{(Combined: \leanref{Substrate.certifier\_silent\_or\_unsound}, module
\leanref{Substrate/Certifier.lean}.)}\par}
\end{theorem}

\begin{proof}[Proof sketch]
Both directions apply \cref{thm:prefix} to the prefix that \(T\) actually inspects
(\(\mathrm{Prefix}(f_0, T.\mathrm{budget})\) for a hypothetical positive verdict \(f_0\)): the
\(\D\)-witness and the \(\mathtt{Cstrict}\)-witness agreeing with that prefix both trigger the
same verdict, so soundness for one class forces the (false) claim that the other class's
witness also lies in the sound class, contradicting \cref{cor:disjoint-dense}.
\end{proof}

A certifier of either substrate class is therefore either unsound (fires on the wrong class at
least once) or permanently silent (never fires at all) --- there is no third option. This is the
model-theoretic counterpart of NEMS Paper~35's programme-internal ``no internal hypercomputer''
result \citep{SpivackP35} (\cref{sec:history}): Paper~35 argues from the
Internality/Outsourcing Schema that no admissible internal procedure may consult an oracle, an
existence claim about hypercomputers; \cref{thm:certifier} complements it with a distinct,
certifier-silence conclusion for this paper's specific model, by the independent route of
prefix-density in Baire space.

\begin{theorem}[Law-Certification Barrier]
\label{thm:lawbarrier}
Let \(f\) be computable, and let \(\mathrm{LawMatches}(f,e)\) (\leanref{Substrate.LawMatches})
say that code \(e\) matches \(f\) at every time step: \(\forall t,\ e.\mathrm{eval}(t) =
\mathrm{Part.some}(f(t))\). Then \(\mathrm{LawMatches}(f, \cdot)\) is not a computable predicate on codes:
\[
  \neg\, \mathrm{ComputablePred}\bigl(\lambda\, e.\ \mathrm{LawMatches}(f,e)\bigr).
\]
{\raggedright\emph{(\leanref{Substrate.law\_match\_undecidable}, module \leanref{Substrate/LawBarrier.lean}.)}\par}
\end{theorem}

\begin{proof}[Proof sketch]
By reduction from the halting problem, reusing the halting-reduction \emph{pattern} of
\texttt{nems-lean}'s \texttt{NemS.Diagonal.HaltingReduction} \citep{nemsLean2025} (no code
dependency; Mathlib only). Given a code \(c\), build a code \(e_c\)
(\leanref{Substrate.matchFn}, \leanref{Substrate.bigG}) that, on input \(t\), outputs \(f(t)\)
unless \(c\) has \emph{visibly} halted on input \(0\) within \(t\) steps (checked via the
step-bounded evaluator \(\mathtt{evaln}\)), in which case it outputs \(f(t) + 1\)
(\leanref{Substrate.matchFn\_eq\_f\_iff}). Then \(e_c\) matches \(f\) forever if and only if
\(c\) never halts on \(0\). The assignment \(c \mapsto e_c\) is computable, uniformly, via the
\(S^m_n\) theorem (\texttt{Nat.Partrec.Code.curry}). If \(\mathrm{LawMatches}(f, \cdot)\) were a
computable predicate on codes, composing with \(c \mapsto e_c\) would computably decide
\(\neg\,\mathrm{HaltsOnZero}(c)\) for every \(c\), hence would decide the halting problem,
contradiction.
\end{proof}

\setcounter{corollary}{0}
\begin{corollary}[Bridge to Paper~30]
\label{cor:p30}
Specializing \cref{thm:lawbarrier} to the constant transcript \(f \equiv 0\) gives
\(\neg\,\mathrm{ComputablePred}(\lambda\,e.\ \mathrm{LawMatches}(0,e))\)
(\leanref{Substrate.law\_match\_undecidable\_zero}). This is the substrate-specific instance of
the general no-total-self-certifier phenomenon proved by NEMS Paper~30, ``Second Incompleteness
for Self-Certifying Learners'' \citep{SpivackP30}: no total internal certifier exists under
diagonal capability. Substrate law-certification is thus a special case of self-certification
failure, not a new phenomenon.
\end{corollary}

\begin{keybox}[title={Interpretation: two closures, not one}]
Theorem~2 says no finite \emph{data} decides substrate. Theorem~4 says no finite \emph{theory}
can be internally certified either, even when the theory is handed to the observer outright.
Together they close both the evidential route and the law-based route; the second closure is
the diagonal barrier proper, and it is the one that genuinely needs computability theory.
\end{keybox}

\begin{thmfourprime}[\(\Pi^0_2\)-hardness of \(\mathrm{LawMatches}\)]
\label{thm:lawbarrier-pi2}
Let \(f\) be computable. The totality problem \(\mathrm{Tot} := \{e \mid \forall t,\
e.\mathrm{eval}(t)\text{ halts}\}\) many-one reduces to \(\mathrm{LawMatches}(f,\cdot)\):
\[
  \mathrm{Tot} \le_m \bigl(\lambda\, e.\ \mathrm{LawMatches}(f,e)\bigr).
\]
Since \(\mathrm{Tot}\) is \(\Pi^0_2\)-complete --- a standard recursion-theoretic fact, cited
here and not proved in this repository --- this pins \(\mathrm{LawMatches}(f,\cdot)\)'s
complexity exactly at \(\Pi^0_2\)-complete, strictly above the co-halting lower bound Theorem~4
gives alone. \emph{(\leanref{Substrate.law\_match\_pi2\_hard}, specialized to \(f \equiv 0\) as
\leanref{Substrate.law\_match\_pi2\_hard\_zero}; module \leanref{Substrate/LawBarrier.lean}.)}
\end{thmfourprime}

\begin{proof}[Proof sketch]
Given a code \(e\), build a code \(e'\) that, on input \(t\), first runs \(e\) on that same
input \(t\) to completion (discarding the result) and then outputs \(f(t)\); this
sequential-composition construction is computable, uniformly in \(e\). If \(e \in \mathrm{Tot}\)
(\(e\) halts on every input), \(e'\) reaches its final step on every \(t\) and outputs \(f(t)\)
there, so \(e'\) matches \(f\) forever; if \(e \notin \mathrm{Tot}\), \(e\) fails to halt on some
particular input \(t_0\), so \(e'\) never reaches its final step on \(t_0\) and fails to match
\(f\) there. This gives \(\mathrm{Tot} \le_m \mathrm{LawMatches}(f,\cdot)\) directly,
independently of Theorem~4's own co-halting reduction.
\end{proof}

% ============================================================================
% 6. WHAT THE THEOREMS DO NOT SAY
% ============================================================================
\section{What the Theorems Do Not Say}
\label{sec:whatnot}

This section exists because the results above are easy to overstate, and an easy target for a
referee who reads past the theorem statements. We state the limits precisely.

\begin{enumerate}[leftmargin=1.8em,itemsep=0.4em]
\item \textbf{The theorems do not say the universe \emph{is} discrete.} They say the question
cannot be settled internally, by evidence (\cref{thm:notfindet}) or by law-certification
(\cref{thm:lawbarrier}), and they identify what \emph{can} adjudicate instead
(\cref{sec:mdl}).

\item \textbf{The theorems do not fully decide axis~C (randomness).} A discrete deterministic
world and a continuum world with objective chance can produce statistically identical Bell
correlations; nothing above distinguishes them. Axis~C is addressed directly by the glimmer
corollary (\cref{sec:glimmer}): its definitions, its density/separation-from-\(\D\) topology, and
its own instance of \cref{thm:notfindet} are machine-checked in a companion Lean module
(\leanref{Substrate/Glimmer.lean}, part of the audited, sorry-free build); the
finite-description consequence of \cref{thm:mdl-uniqueness} depends on this class existing at
all, so \cref{sec:glimmer} states it as settled content rather than an open problem. What remains genuinely open,
independently of that Lean status, is the single \textbf{Born-martingale lemma} --- that the
actual physical Born measure, as complete internal semantics per NEMS Papers~13/14
\citep{SpivackP13,SpivackP14}, admits no successful computable martingale on the realized
transcript, given the non-degeneracy premise of NEMS Paper~12 \citep{SpivackP12} and the degree
bound of UGP Paper~51 \citep{SpivackUGP51}. We state the dependency chain honestly in
\cref{sec:glimmer} and do not claim more of the corollary than is actually established: the
mathematical machinery (the definitions, the separations, and the elementary bet-everything
martingale fact) is machine-checked; the physical completeness premise connecting it to the
actual Born measure is not.

\item \textbf{The theorems do not say discreteness is unfalsifiable in every sense.}
\cref{thm:notfindet} rules out \emph{finitely determined} (clopen) properties, but asymptotic
properties --- limits over unboundedly long observation, not settled by any fixed prefix
length --- are not excluded and are in principle limit-testable (\cref{sec:scaling}). The
theorem locates the only place substrate information can live; it does not eliminate all
empirical content.

\item \textbf{The continuum class in this v1 development is classical and unrestricted}, so
\(\mathtt{C}\) is all of Baire space (\cref{rem:modeling}). A computable-analysis refinement
would shrink \(\mathtt{Cstrict}\), not empty it, by Pour-El--Richards
\citep{PourElRichards1981,PourElRichards1989}: their non-computable wave-equation solution
already arises from computable initial data, so even a maximally restrictive computable-analysis
reading of ``continuum'' still admits non-computable transcripts. The results are robust to
this refinement; we have not carried it out because Mathlib does not yet provide computable
analysis (\cref{sec:open}).

\item \textbf{Theorem~2 does not depend on Church--Turing; Theorem~4 does, and the oracle
strength needed to escape it is pinned down exactly.} \cref{thm:notfindet} is pure density
in Baire space (the boxed remark after its proof) and would hold verbatim for any two dense,
disjoint transcript classes, computable or not; no oracle, at any level, changes which prefixes
exist in \(\D\) and \(\mathtt{Cstrict}\) or the density fact between them, so \cref{thm:notfindet}
is not escaped by \emph{any} amount of extra computational power. \cref{thm:lawbarrier} is a
genuine diagonal argument that requires the observer's internal test to be a computable
predicate --- and the oracle strength that defeats it is exactly \(\mathbf{0}''\), not
\(\mathbf{0}'\). \(\mathrm{LawMatches}(f,\cdot)\) is a \(\Pi^0_2\) predicate on codes (a
universal claim, over every time step, that evaluation halts \emph{and} agrees with \(f\)); the
reduction proving \cref{thm:lawbarrier} gives co-halting-hardness, and a second reduction
(\leanref{Substrate.law\_match\_pi2\_hard}: the totality problem \(\mathrm{Tot} := \{e \mid
\forall t,\ e.\mathrm{eval}(t) \text{ halts}\}\) many-one reduces to \(\mathrm{LawMatches}(f,\cdot)\),
via a code that runs the candidate on each input to completion and discards the result)
establishes \(\Pi^0_2\)-\emph{hardness} --- strictly above the co-halting lower bound alone,
and, since \(\mathrm{Tot}\) is \(\Pi^0_2\)-complete (a standard, cited recursion-theoretic fact,
not proved in this repository), this pins \(\mathrm{LawMatches}(f,\cdot)\)'s complexity exactly
at \(\Pi^0_2\)-complete. A \(\mathbf{0}''\) oracle (halting relative to \(\mathbf{0}'\),
sufficient to decide any \(\Pi^0_2\) statement) escapes \cref{thm:lawbarrier}; a plain halting
(\(\mathbf{0}'\)) oracle provably does \emph{not}, since no \(\Pi^0_2\)-complete problem is
decidable relative to \(\mathbf{0}'\) --- superseding the earlier, weaker claim that a
\(\mathbf{0}'\) oracle was merely ``not shown'' to escape. What \emph{is} established, and is
worth stating explicitly, is that \cref{thm:notfindet} survives strictly more observer upgrades
than \cref{thm:lawbarrier}: no oracle escapes the former, while exactly \(\mathbf{0}''\)-and-above
escapes the latter.
\end{enumerate}

% ============================================================================
% 7. WHERE SUBSTRATE INFORMATION CAN LIVE: SCALING EXPONENTS
% ============================================================================
\section{Where Substrate Information Can Live: Scaling Exponents}
\label{sec:scaling}

\cref{thm:notfindet} rules out finitely determined (clopen) observables as substrate evidence.
It does not rule out \emph{asymptotic} observables --- quantities defined as a limit over an
unboundedly growing resource budget, which by construction are not settled by any fixed prefix
length. This section classifies the known candidates (status: discussion-level, ``D'' in
Appendix~\ref{app:status}, not machine-checked) and explains, for each, why it is the correct
\emph{kind} of empirical target under the barriers proved above.

\paragraph{Metrological scaling.} In an unrestricted continuum, achievable measurement precision
scales without a floor as the number of resources (e.g.\ repeated trials or entangled probes)
\(N\) grows --- the Heisenberg \(1/N\) scaling law. Under finite substrate information, this
scaling exponent must break at some resource level, because there are only finitely many
distinguishable outcomes to resolve into. The signature is an exponent measured over decades of
resource scale, not a sharp threshold at any single \(N\) --- which is exactly the form
\cref{thm:notfindet} allows. This signature is forced by discreteness of \emph{any} kind,
unlike a Lorentz-violation search, which tests only one specific discretization scheme
(\cref{sec:history}).

\paragraph{Spectral type and recurrence.} A finite-dimensional Hilbert space forces pure point
spectrum, exact quasi-periodicity, and a finite Poincar\'e recurrence time; a continuum Hilbert
space permits continuous spectra and, generically, no finite recurrence time. This is
observable, in principle, as a recurrence-time scaling law with system size, again an
asymptotic rather than a threshold signature.

\paragraph{Complexity class of physical computation (axis~B\('\), not axis~B).} Any physical
process solving a problem outside \(\mathsf{BQP}\) with polynomial physical resources would be
evidence on axis~B\('\) --- efficient computability (\cref{sec:intro}, \cref{tab:axes}) --- not
axis~B: such a process remains computable in the ordinary sense, so it says nothing about
membership in \(\mathtt{Cstrict}\) and does not itself bear on \cref{thm:lawbarrier}. It is
nonetheless the natural empirical target for Aaronson's programme \citep{Aaronson2005}
(\cref{sec:history}). By the same prefix-density argument (\cref{thm:prefix}) underlying
\cref{thm:certifier}: a \emph{single} such demonstration can never be internally certified as a
genuine complexity-class violation, since any finite observed prefix of the demonstration is
equally consistent with an ordinary, bound-respecting computable continuation, so no internal
test built from that prefix could be sound and ever fire. Only the \emph{scaling} of resource
against instance size, across many instances, can be reported as evidence --- never certified
as proof.

\paragraph{Constants.} Finite decimal expansions of physical constants decide nothing:
\cref{thm:notfindet} says any finite data is compatible with both \(\D\) and \(\mathtt{Cstrict}\).
A theory that predicts a constant from a finite combinatorial structure and keeps being
empirically confirmed to increasing precision is MDL \emph{evidence} in the sense of
\cref{thm:outsourcing} --- never certification in the sense of \cref{thm:lawbarrier}. This is
explicitly the correct epistemic status to claim for predictions coming out of the Universal
Generative Principle (UGP) programme, and we state it plainly: such predictions, however
striking, are evidence of the kind Theorem~5 identifies as the unique admissible kind, not
proof of the kind Theorem~4 shows to be structurally unavailable.

% ============================================================================
% 8. THEOREM 5: OUTSOURCING AND MDL
% ============================================================================
\section{Theorems, Part II: Outsourcing and the Unique Admissible Adjudicator}
\label{sec:mdl}

\begin{definition}[Internally decidable]
\label{def:intdec}
A transcript property \(P\) is \emph{internally decidable}
(\leanref{Substrate.InternallyDecidable}) if some internal test decides it exactly:
\(\exists\, T,\ \forall f,\ T.\mathrm{run}(f) = \mathrm{true} \leftrightarrow P(f)\).
\end{definition}

\begin{theorem}[Outsourcing and the Unique Adjudicator]
\label{thm:outsourcing}
\label{thm:mdl-uniqueness}
\begin{enumerate}[label=(\alph*), leftmargin=1.6em]
\item \textbf{Outsourcing (Lean-checkable half).} No internally decidable property separates
\(\D\) from \(\mathtt{Cstrict}\):
\[
  \neg\, \exists\, P,\ \mathrm{InternallyDecidable}(P) \wedge (\forall f \in \D,\ P(f)) \wedge
  (\forall f \in \mathtt{Cstrict},\ \neg P(f)),
\]
and symmetrically for the reverse separation
(\leanref{Substrate.substrate\_not\_internally\_decidable},
\leanref{Substrate.substrate\_not\_internally\_decidable'}). Consequently, neither \(f \in \D\)
nor \(f \in \mathtt{Cstrict}\) is itself internally decidable
(\leanref{Substrate.discrete\_not\_internally\_decidable},
\leanref{Substrate.strict\_continuum\_not\_internally\_decidable}). \emph{(Module
\leanref{Substrate/Outsourced.lean}; status \A.)}
\item \textbf{Uniqueness of the MDL adjudicator (cited, not re-derived).} The only adjudication
procedure consistent with PSC --- one that (i) consults no oracle and (ii) is invariant under
presentation of the same law-set --- is selection over complete law-sets by minimum description
length (MDL). \emph{(Argument cited to NEMS Paper~34, ``A Sieve Engine for Theory Spaces,'' and
Paper~84, ``The Survivor Calculus'' \citep{SpivackP34,SpivackP84}; status \AMDL, \emph{not} \A{}
--- this half is not a Mathlib reduction and is not re-derived here.)}
\end{enumerate}
\end{theorem}

\begin{proof}[Proof sketch, part (a)]
Every internally decidable property is finitely determined
(\leanref{Substrate.internally\_decidable\_implies\_finitely\_determined}: a test only ever
inspects a bounded prefix), so part (a) follows immediately from \cref{thm:notfindet}. Part (b)
is cited to the NEMS programme results named above, not re-derived in this repository.
\end{proof}

Under Perfect Self-Containment (PSC), a self-contained system has no external selector to
consult: there is no oracle, and no observer outside the system, that could supply a verdict the
system itself cannot compute. \cref{thm:notfindet,thm:lawbarrier} together show that substrate
class is not internally decidable from evidence or from law-certification --- part (a) above.
Hence substrate discreteness is an \emph{outsourced} property in the precise sense of the
Internality/Outsourcing Schema, NEMS Paper~83 \citep{SpivackP83}: a load-bearing task (here,
adjudicating substrate) that cannot be discharged by any procedure internal to the system doing
the adjudicating --- which is exactly what makes part (b)'s external adjudicator necessary.

Consequence: \(\D\) consists exactly of the transcripts admitting a finite description (by
\cref{lem:transcript-char}, a computable one); \(\mathtt{Cstrict}\) by construction admits none.
MDL therefore adjudicates, uniquely and non-arbitrarily, in favor of the \textbf{finite-description
class (\(\D \cup \mathrm{Glimmer}\))}, not \(\D\) alone --- not as a methodological preference for
simplicity, but as the one selection procedure that survives the closure requirement PSC
imposes. \(\D\) alone would be the wrong name for what MDL selects: NEMS Paper~12 (Determinism
No-Go) proves that no total computable law generates the records in a diagonal-capable PSC
universe with genuine record-divergent choice \citep{SpivackP12}, so a universe with that
property has no computable transcript and hence is not itself a member of \(\D\); the glimmer
class (\(\Delta^0_2\), finitely described by law plus one jump, \cref{sec:glimmer}) is what
closes this gap, and is why \cref{sec:glimmer} is a load-bearing part of this theorem's
statement rather than a deferred open problem. This is the bridge to the Universal
Generative Principle (UGP) programme: UGP's stance that a compact, finite generative structure
functions as a certificate for physical law is precisely \cref{thm:mdl-uniqueness} applied to
one specific law-set, and \cref{sec:scaling}'s remark on constants states the resulting
epistemic status explicitly.

\paragraph{``Discrete'' is a description-length notion, not a geometry, and a continuum
substrate can be its worked example.} One clarification matters for a physicist reader applying
\cref{thm:mdl-uniqueness} to a candidate fundamental theory: \emph{``discrete'' in this paper
never meant ``lattice''; it meant ``finitely generated.''} \(\D\) is defined by computable
\emph{transcript} (\cref{lem:transcript-char}), a description-length class, not a geometry
class, and \cref{lem:embedding} already proves \(\D \subseteq \mathtt{C}\): every
discrete-transcript world is realized by \emph{some} continuum world. A continuum substrate with
a finite arithmetic generator therefore sits in \(\D \cap \mathtt{C}\) by transcript ---
continuum by state type, finitely described by transcript --- and this is not a counterexample
to \cref{thm:mdl-uniqueness}; it is the theorem's worked example. The Universal Generative
Principle (UGP) programme's \(\Phi_{\mathrm{MDL}}\)/GTE framework is exactly this case: a
continuum field whose observable transcript is generated by a finite arithmetic certificate, so
it is precisely the kind of world \cref{thm:mdl-uniqueness} selects, stated in the vocabulary of
a specific physical theory rather than in the abstract model above --- not a rival to this
paper's agnosticism, but an instance of what the theorem already predicts should be favored.

This also resolves an apparent tension with NEMS Paper~24's claim that the universe must begin
as a discrete, self-interpreting ``Reflexive Seed'' \citep{SpivackP24}. That claim is a
\cref{thm:mdl-uniqueness}-direction result about \emph{initial conditions}, not an internal
certification of the \cref{thm:lawbarrier} kind: a classical singularity is an underspecified
boundary that requires external completion, hence non-foundational under PSC, so closure
prefers a finite, self-interpreting origin over an unbounded one \citep{SpivackP18,SpivackP27}.
Crucially, the seed is prior in \emph{generative} order, not \emph{temporal} order. The
Infinity Compression thesis's two-layer structure --- a bare certificate \(B\), an enriched
realization \(E\), and a comparison map \(\pi : E \to B\) with nontrivial fibers and sections
--- is the formal handle for exactly this relation: the finite generator is \(B\), the continuum
realization is \(E\), and asking whether \(B\) exists ``before'' \(E\) is asking whether a
genotype exists before its phenotype (an analogy offered with its own caveat, since it is not
exact: a genotype is itself a physical object in time, whereas the seed is a type) --- generative
priority, not temporal priority. On this reading, Paper~24's discreteness claim and this paper's
substrate-agnosticism are consistent rather than in tension: Paper~24 makes a claim about
admissible initial data (an outsourcing/MDL selection, in this paper's \cref{thm:mdl-uniqueness}
sense), not a claim that spacetime itself is a lattice, and not an internal certification that
would run afoul of \cref{thm:lawbarrier}.

\setcounter{corollary}{0}
\begin{corollary}[Uncountable-observer strengthening of Theorem~5]
\label{cor:cardinality}
\(\D\) is countable (\leanref{Substrate.D\_countable}) and \(\mathtt{Cstrict}\) is not
(\leanref{Substrate.Cstrict\_not\_countable}). More strongly, no family of internal tests
\(\mathrm{tests} : I \to \mathrm{InternalTest}\), indexed by an \emph{arbitrary} type \(I\) ---
including an uncountable one --- jointly decides discrete membership:
\begin{multline*}
  \neg\, \exists\, (\mathrm{tests} : I \to \mathrm{InternalTest}),\
  (\forall f \in \D,\ \exists i,\ \mathrm{tests}(i).\mathrm{run}(f)) \,\wedge\, {} \\
  (\forall f \in \mathtt{Cstrict},\ \forall i,\ \neg\,\mathrm{tests}(i).\mathrm{run}(f)).
\end{multline*}
{\raggedright\emph{(\leanref{Substrate.no\_joint\_substrate\_decision}, module
\leanref{Substrate/Cardinality.lean}; an analogous strengthening for internally decidable
predicates, over any index type, is \leanref{Substrate.no\_joint\_internally\_decidable}.)}\par}
\end{corollary}

\begin{proof}[Proof sketch]
The obstruction is exactly \cref{thm:prefix} again, not a cardinality argument: fix any
candidate family and let \(f_0\) be the constant-zero (discrete) transcript; some test \(i\)
must accept \(f_0\). Applying \cref{thm:prefix} to the prefix that test \(i\) actually inspects
produces a \(\mathtt{Cstrict}\)-witness agreeing with \(f_0\) on that prefix, so test \(i\)
accepts it too, contradicting the joint-decision hypothesis. No step of this argument depends
on the cardinality of the index type \(I\); an uncountable family of observers adds no power
beyond a single one (\leanref{Substrate.uncountable\_index\_offers\_no\_power}). The
countability facts above are worth recording in their own right --- \(\D\) is a subset of the
computable functions, hence countable (\leanref{Substrate.computable\_set\_countable}), while
\(\mathtt{Cstrict}\) occupies almost all of Baire space (\leanref{Substrate.baire\_space\_not\_countable})
--- but they are not what does the work here; the work is prefix indistinguishability, exactly
as in \cref{thm:notfindet,thm:outsourcing}.
\end{proof}

This forecloses a natural escape attempt one might otherwise propose against \cref{thm:outsourcing}:
pooling an unlimited, even uncountable, supply of cooperating internal observers cannot recover
what a single observer cannot decide, since the obstruction is topological density in Baire
space, which is entirely insensitive to how many observers are asking.

\begin{corollary}[Measure-theoretic typicality of the continuum class]
\label{cor:measure-typicality}
Under the canonical geometric i.i.d.\ product measure \(\mu_p\) on Baire space (any parameter
\(p \in (0,1)\)), \(\D\) has measure zero and \(\mathtt{Cstrict}\) has full measure:
\[
  \mu_p(\D) = 0, \qquad \mu_p(\mathtt{Cstrict}) = 1.
\]
{\raggedright\emph{(\leanref{Substrate.D\_null\_geometric}, \leanref{Substrate.Cstrict\_full\_mass\_geometric},
module \leanref{Substrate/Probabilistic.lean}.)}\par}
\end{corollary}

\begin{proof}[Proof sketch]
\(\mu_p\) has no atoms: the per-coordinate geometric measure of any single symbol is bounded
above by \(p < 1\) uniformly, so the infinite product of these bounds tends to \(0\), giving
\(\mu_p(\{f\}) = 0\) for every transcript \(f\). Since \(\D\) is countable (\cref{cor:cardinality}),
it is a countable union of \(\mu_p\)-null singletons, hence itself \(\mu_p\)-null
(\leanref{Substrate.D\_null\_geometric}). Since \(\mathtt{C} = \mathrm{univ}\)
(\cref{rem:modeling}, \leanref{Substrate.C\_eq\_univ}), \(\mathtt{Cstrict} = \mathrm{univ}
\setminus \D\) (\leanref{Substrate.Cstrict\_eq\_univ\_diff\_D}), so
\(\mu_p(\mathtt{Cstrict}) = 1 - \mu_p(\D) = 1\)
(\leanref{Substrate.Cstrict\_full\_mass\_geometric}).
\end{proof}

This reinforces \cref{cor:cardinality} from a different direction: not only does no family of
internal observers, of any cardinality, jointly decide substrate class, but under the geometric
family of product measures specifically, a transcript drawn at random is strict-continuum with
probability~1. That verdict is not, however, forced by the mathematics of measure itself, and
stating it as typicality \emph{simpliciter} would be exactly the external-selector move
\cref{sec:psc} identifies as PSC's target: choosing \(\mu_p\) over every other measure on Baire
space is itself an unexplained external choice, on a par with eternal inflation's measure
problem. The clearest way to see this is to change the choice. The Solomonoff/MDL universal
prior --- \cref{sec:psc}'s own example of PSC applied to constants, and the actual adjudication
principle \cref{thm:mdl-uniqueness} invokes --- weights a sequence's finite prefixes by
description length rather than by a fixed per-symbol probability: a computable sequence
generated by a program of size \(c\) has prefix-length-\(n\) algorithmic probability bounded
below by \(2^{-c - O(\log n)}\), a decay that is polynomial in \(n\), while \emph{every}
sequence's length-\(n\) prefix, computable or not, decays exponentially in \(n\) under any
nondegenerate \(\mu_p\) --- the identical atomlessness mechanism behind this corollary's own
proof \citep{LiVitanyi2008}. The universal prior itself depends on a choice of reference
machine, but only up to an additive constant fixed by the pair of machines being compared, not
by the sequence under evaluation --- the invariance theorem's own guarantee
\citep{LiVitanyi2008} --- which is exactly presentation-invariance, \cref{thm:mdl-uniqueness}'s
condition~(ii), applied to the choice of universal prior itself rather than to a choice among
law-sets: the universal prior is therefore an admissible ingredient for the kind of
closure-respecting adjudication \cref{thm:mdl-uniqueness} already requires, not a second,
unexamined external choice smuggled in alongside \(\mu_p\)'s. Under the universal prior, then,
finitely describable sequences are
the ones exponentially favored, not the ones typicality discards: \cref{cor:measure-typicality}'s
verdict and the universal prior's verdict are exact opposites, and the entire disagreement
between them lives in which prior the adjudicator adopts, not in any fact about the substrate.
This is the sharpest illustration available of why \cref{thm:mdl-uniqueness} adjudicates the way
it does: it does not sample from \(\mu_p\), from the universal prior, or from any measure over
Baire space at all, precisely because doing so would smuggle back in the external selector PSC
forbids. MDL selection over complete law-sets is a closure-respecting procedure, not a draw from
a distribution over which universe is real, and \cref{cor:measure-typicality} is best read not as
evidence for the continuum but as a demonstration of exactly how badly a naive sampling-measure
adjudicator fails the very closure requirement this paper's actual adjudicator is built to
satisfy. None of this bears on what an \emph{internal test} can detect from a finite prefix, nor
does it by itself yield an \(\varepsilon\)-tolerant test: see the probabilistic-observers
discussion in \cref{sec:open} for exactly where that further step remains open and why.

\subsection{The glimmer corollary}
\label{sec:glimmer}

The glimmer corollary is essential to Theorem~5, not an optional addendum to it: without it,
the theorem's conclusion would name a class that the programme's own account of axis~C says
the actual universe is not in. NEMS Paper~12
(Determinism No-Go) proves that no total computable law generates the records in a
diagonal-capable PSC universe with genuine record-divergent choice \citep{SpivackP12}, so a
universe with that property does \emph{not} have a computable transcript, and \(\D\) alone would
then be the wrong name for what MDL actually selects. Glimmer resolves this: a glimmer world has
a \emph{computable law} and a \(\Delta^0_2\) transcript --- the outcomes are not free bits, since
they are fixed by the law plus one jump with no external selector, so the world still has a
finite description and \cref{thm:mdl-uniqueness} still selects it. \cref{sec:mdl}'s consequence
paragraph and the abstract accordingly say ``finite-description class (\(\D \cup
\mathrm{Glimmer}\))'' rather than ``finite-information class,'' for exactly this reason. The
definitions and theorem statements below are proved in a companion Lean module,
\leanref{Substrate/Glimmer.lean}, \textbf{machine-checked and part of this repository's audited,
sorry-free build} (full Lean status disclosed below, after the derivation chain); the
classification itself is a matter of mathematics and formalization, complete and
machine-checked --- only the Born-martingale lemma below carries an open physical premise.

\emph{Informal idea.} A \emph{glimmer} is a sequence that is random from inside an effective
description and structured from one level above it. The paradigm example is Chaitin's
\(\Omega\) \citep{Chaitin1975}: no computable betting strategy (martingale) beats it, yet it is
fixed by a finite formula and is computable \emph{from} the halting set --- it sits at Turing
degree \(\mathbf{0}'\), i.e.\ it is \(\Delta^0_2\). (\(\Omega\) is in fact Martin-L\"of random,
a strictly stronger property than immunity to merely computable martingales; the formal
definition below deliberately uses the weaker \emph{computably random} notion, since that is
exactly what the Born-martingale lemma's argument, stated in terms of computable martingales, is
capable of establishing --- see the derivation chain below. Since Martin-L\"of randomness
implies computable randomness, \(\Omega\) satisfies the weaker notion too, so the flagship
example is unaffected by this choice.) The claim under investigation, not proved here, is that
quantum measurement/adjudication outcomes are glimmers relative to the physical law: not brute
randomness, but structure the law itself cannot parse from the inside.

\emph{Idealization caveat.} Randomness notions --- computable randomness, Martin-L\"of randomness
--- are properties of \emph{infinite} sequences. Nothing here commits the paper to an infinite
universe: in a finite-information, finite-lifetime universe every realized transcript is finite,
and the correct finitary form of the claim below is \emph{prefix incompressibility}, \(K(f
\!\restriction\! n) \ge n - c\) for the realized prefixes relative to the law, for a constant
\(c\) independent of \(n\). We state the infinite-sequence form throughout for continuity with
the algorithmic-randomness literature cited below; the finitary form is the one that actually
applies to any physically realized, finite-lifetime transcript.

\emph{Four-class taxonomy.} Fix an effective description \(E\) = (computable law with
computable Born measure \(\mu\), internal observer) and let \(f : \N \to \N\) be the
adjudication transcript. The four classes below exhaust all possible adjudication transcripts
relative to \(E\) precisely because UGP Paper~51's bound \citep{SpivackUGP51} establishes
\(f \le_T \mathbf{0}'\) under that paper's hypotheses (the DSAC model) for transcripts arising in
this framework, ruling out a residual fifth case --- non-computable, not computably random, and
not \(\Delta^0_2\) --- that would otherwise fall outside all four boxes below.

\begin{center}
\small
\begin{tabular}{@{}>{\raggedright\arraybackslash}p{2.4cm}>{\raggedright\arraybackslash}p{5.2cm}>{\raggedright\arraybackslash}p{2.4cm}>{\raggedright\arraybackslash}p{4.2cm}@{}}
\toprule
\textbf{Class} & \textbf{Definition} & \textbf{Example} & \textbf{Status} \\
\midrule
1. Pseudorandom & \(f\) computable & any PRNG & excluded by the non-degeneracy premise \citep{SpivackP12} plus the Born-martingale lemma (below); \emph{not} by Paper~11 alone \\
2\('\). Unpredictable, compressible & \(f\) not computable, not computably random, \(f \le_T \mathbf{0}'\) & the halting set itself & excluded by the \emph{same} lemma as class~1 --- see the derivation chain below \\
2. \textbf{Glimmer} & \(f\) computably random relative to \(\mu\) \textbf{and} \(f \le_T \mathbf{0}'\) (i.e.\ \(\Delta^0_2\)) & \(\Omega\) & the claimed class for physical randomness \\
3. Truly random & \(f\) computably random relative to \(\mu\), not \(\Delta^0_2\) & a.e.\ Martin-L\"of-random sequence & excluded by UGP Paper~51 \citep{SpivackUGP51} \\
\bottomrule
\end{tabular}
\end{center}

Class~2 is measure-zero but nonempty. The distinction between \emph{unpredictable} (no total
computable predictor) and \emph{incompressible} (immune to every computable martingale) is
exactly what separates class~1 from class~2\('\); NEMS Paper~11 (``Physical Incompleteness''
\citep{SpivackP11}) is cited only as context for this distinction, not as doing the
class-1-exclusion work itself --- that work is done by the derivation chain that follows.

\emph{Derivation chain.} Class~1 and class~2\('\) exclusion does not route through Paper~11, nor
through any premise about the realized transcript being non-computable as a \emph{conclusion} of
a law-level result; instead it falls out of the
Born-martingale lemma itself, with Paper~12 supplying a premise rather than a conclusion.
\begin{enumerate}[leftmargin=1.8em,itemsep=0.3em]
\item \textbf{Premise (non-degeneracy / genuine choice).} \(\mu(\text{next symbol}) \le 1 -
\varepsilon\) infinitely often along \(f\), for some fixed \(\varepsilon > 0\). This is exactly
NEMS Paper~12's own premise, ``record-divergent choice'' \citep{SpivackP12}, so Paper~12 is cited
\emph{for the premise}, not for a claim about the transcript's computability. A degenerate
\(\mu\) (all conditional probabilities eventually \(0\) or \(1\)) gives no randomness and the
question does not arise.
\item \textbf{Born-martingale lemma (open).} If \(\mu\) is complete internal semantics per
Papers~13/14 \citep{SpivackP13,SpivackP14}, no computable \(\mu\)-martingale succeeds on \(f\)
--- i.e.\ \(f\) is computably random relative to \(\mu\). Combined with the elementary fact that
a computable, non-degenerate \(f\) admits a computable martingale that bets everything on
\(f(t)\) and succeeds (\leanref{Substrate.computable\_not\_random}, module
\leanref{Substrate/Glimmer.lean}), this gives \(f \notin\) class~1 and \(f \notin\) class~2\('\)
simultaneously, since both classes require \(f\) to fail computable randomness. \textbf{The
completeness premise itself --- that the physical Born measure genuinely constitutes complete
internal semantics, hence that no computable martingale succeeds on the actually realized
transcript --- is not proved anywhere in the programme}; the gap is narrow and precisely
located, not generic. It is \emph{not} a randomness-theory gap: the martingale and
computable-randomness machinery the lemma needs is already fully built and machine-checked
(\cref{sec:glimmer}). Nor can it be closed by simply \emph{defining} ``complete internal
semantics'' to mean ``\(f\) is computably random'' --- that would presuppose exactly what the
lemma is meant to establish, foreclosed by this development's no-smuggling convention; indeed
the unrestricted universal form of the premise (every non-degenerate \(f\) admits no successful
computable martingale) is outright \emph{false} here, since
\leanref{Substrate.BetMartingale.bet\_succeeds} proves the opposite for every \emph{computable}
non-degenerate \(f\), so the lemma can only ever concern the specific, presumptively
non-computable, realized transcript of a genuinely complete-semantics universe. Nor is a
suitable bridge available from the existing quantum-formalism side: the nearest existing
candidate for ``complete internal semantics,'' \texttt{BICS} in the companion
\texttt{nems-lean} development (Papers~13/14), is a semantic-architecture predicate over
density operators and Born-trace probabilities, with no notion of a realized \(\N \to
\mathrm{Fin}\,k\) sequence or computable martingale at all --- categorically the wrong kind of
object to bridge directly. What remains open is therefore exactly and only a
non-question-begging formalization of the physical premise itself as a predicate on realized
transcripts; this is the one piece that cannot be closed without new physical or interpretive
input, and this paper states the dependency honestly rather than treating it as established.
\item \textbf{UGP Paper~51 (adjudication at degree exactly \(\mathbf{0}'\)).} \(f \le_T
\mathbf{0}'\) \citep{SpivackUGP51} \(\Rightarrow\) \(f \notin\) class~3.
\item Hence, conditionally on step~2's premise, \(f \in\) class~2 (glimmer).
\end{enumerate}
We note explicitly that step~2's supporting fact
(\leanref{Substrate.computable\_not\_random}) establishes only computable-randomness immunity
(against computable martingales), not the strictly stronger Martin-L\"of randomness (which
requires immunity to the larger class of left-c.e.\ lower-semicomputable supermartingales); this
is exactly why the taxonomy above is stated in terms of computable randomness rather than
Martin-L\"of randomness --- using the stronger notion here would require a correspondingly
stronger, less obviously justified premise about what the Born measure's completeness rules out.

\begin{quote}
\small\raggedright
\textbf{Lean status --- \A, machine-checked.} A companion Lean module,
\leanref{Substrate/Glimmer.lean}, proves exactly the definitions and theorem statements named
above (\leanref{Substrate.LimitComputable}, \leanref{Substrate.BornMeasure},
\leanref{Substrate.BornMeasure.NonDegenerateAlong}, \leanref{Substrate.BornMeasure.CMartingale},
\leanref{Substrate.Succeeds}, \leanref{Substrate.ComputablyRandom},
\leanref{Substrate.Glimmer}, \leanref{Substrate.limitComputable\_dense},
\leanref{Substrate.notLimitComputable\_dense},
\leanref{Substrate.limitComputable\_not\_finitely\_determined},
\leanref{Substrate.glimmer\_finite\_description}, \leanref{Substrate.computable\_not\_random},
\leanref{Substrate.glimmer\_not\_computable}, \leanref{Substrate.glimmer\_disjoint\_D},
\leanref{Substrate.glimmer\_subset\_Cstrict}), and is part of this repository's audited,
sorry-free build (imported by \leanref{Substrate.lean}; \texttt{lake build}: 8262/8262 jobs,
zero errors). The bet-everything martingale construction
(\leanref{Substrate.BetMartingale.bet} / \leanref{Substrate.BetMartingale.bet\_succeeds}) is
fully proved, with no \texttt{sorry} placeholders, via an explicit ``wrong-child'' dead-branch
mechanism handling off-path martingale fairness; \leanref{Substrate.glimmer\_finite\_description}
states the non-vacuous fact that a limit-computable sequence is fixed by one code plus a
convergent stage function (\(\exists\, c,\ g,\ \mathrm{Computable}(g) \wedge (\forall t,\
\exists s_0,\ \forall s \ge s_0,\ g(t,s) = f(t)) \wedge \mathrm{eval}(c) = \ldots\)), not merely
an unconditional witness. Every declaration named above is confirmed by direct
\texttt{\#print axioms} to depend on at most the three standard Mathlib axioms
(\texttt{propext}, \texttt{Classical.choice}, \texttt{Quot.sound}); see Appendix~\ref{app:lean}
for the full table. Non-emptiness of \(\mathrm{Glimmer}\) for the uniform measure (Chaitin's
\(\Omega\)) is cited, not built, per Downey--Hirschfeldt \citep{DowneyHirschfeldt2010} ---
status \A/\textsc{d}, unaffected by the point above.
\end{quote}

\emph{Modeling note.} The \texttt{World} structure (\cref{def:world}) is deterministic, so a
glimmer world is not itself a \texttt{World}: this paper's transcripts are the observer's
channel, and the formal model is deliberately class-agnostic about how the substrate produced
them. A glimmer transcript is one \emph{produced by} a deterministic law plus a \(\Delta^0_2\)
adjudication stream, and \(\mathrm{Glimmer}\) is defined directly on transcripts, not on
\texttt{World}. Generalizing \texttt{World} with an explicit adjudication input is future work;
nothing in Part~I needs it.

\emph{What it is not.} Not a hidden-variable theory: the \(\Delta^0_2\) description is of the
whole sequence, not a local pre-assignment, so Bell's theorem \citep{Bell1964} is untouched. Not
internally usable determinism: \cref{thm:notfindet} still applies (class~2 versus class~3 is
itself not finitely determined) and \cref{thm:certifier} still applies (no internal certifier
can fire on this distinction either). The bridge to physical adjudication traces, if the
Born-martingale lemma is eventually proved, would be follow-up work via NEMS
Papers~76/77 \citep{SpivackP76,SpivackP77}.

\emph{Usage rule for this paper and its companions.} ``Glimmer'' means class~2 only, always
relative to a named \(E\). Unqualified ``random'' means classes \(2 \cup 3\); ``truly random''
means class~3 only; ``unpredictable'' means \(\notin\) class~1 and must not be conflated with
``random.''


% ============================================================================
% 9. DISCUSSION
% ============================================================================
\section{Discussion}
\label{sec:discussion}

\subsection{Methodological consequence: hunt exponents, not thresholds}

The central methodological claim of this paper is that ``test whether the substrate is
discrete'' is not merely a hard experiment to design --- it is, for any test of the finitely
determined shape, structurally the wrong kind of question, because \cref{thm:notfindet} proves
no such test can in principle succeed, independent of instrumental precision or ingenuity. This
reframes what counts as progress. A negative result from a Planck-scale interferometry search,
or from a search for Lorentz-violating dispersion, is not evidence \emph{against} discreteness
in the way it would naively be read; it is evidence against \emph{one particular
discretization scheme}, exactly as the causal-set programme demonstrated
(\cref{sec:history}). The correct empirical strategy, per \cref{sec:scaling}, is to hunt
\emph{asymptotic scaling exponents} --- metrological precision scaling, recurrence-time scaling,
resource-versus-instance-size scaling for candidate hypercomputational demonstrations --- because
these are the one class of observable that is not finitely determined and therefore is not
excluded by \cref{thm:notfindet}. This is a genuine redirection of experimental priority, not a
weakening of empirical ambition: it says the interesting physics is in the exponent, not in the
threshold.

\subsection{What is proved versus what is cited: the A\_Lean / A\_MDL line}

It matters, and we state it here without hedging, exactly where the machine-checked boundary
sits. Lemmas~2.1--2.2, Theorems~1--4, Corollary~4.1, Theorem~4\('\)'s \(\Pi^0_2\)-hardness
sharpening (\leanref{Substrate.law\_match\_pi2\_hard}), the outsourcing half of Theorem~5
(\cref{thm:outsourcing}), its uncountable-observer strengthening (\cref{cor:cardinality}), the
deterministic half of the probabilistic-observers result
(\leanref{Substrate.acceptance\_on\_D\_iff\_Cstrict}), the measure-zero/full-measure
classification of \(\D\)/\(\mathtt{Cstrict}\) (\cref{cor:measure-typicality}), the glimmer
corollary's definitions and separations (\cref{sec:glimmer}), and the axis-A dense pair
(\cref{rem:axisA}) are all proved and machine-checked (\A, Appendix~\ref{app:status}) in Lean~4
against Mathlib, with zero \texttt{sorry} and zero custom axioms. The only exceptions on the
Lean side of the ledger are the Born-martingale lemma's physical completeness premise
(\cref{sec:glimmer}, genuinely open, not a formalization gap) and the \(\varepsilon\)-tolerant
randomized generalization of the probabilistic-observers result (\cref{sec:open}). The
\emph{uniqueness}-of-MDL half of Theorem~5 (\cref{thm:mdl-uniqueness}) is \AMDL: an argument
cited to NEMS Papers~34 and~84 \citep{SpivackP34,SpivackP84}, not a Mathlib reduction carried
out in this repository. This is not a gap we are trying to paper over --- it reflects a genuine
difference in the kind of claim being made. ``No internally decidable property separates the
two transcript classes'' is a closed mathematical statement about a formally defined class of
procedures, and Lean can check it end to end. ``MDL is the unique selection procedure meeting
two stated closure desiderata'' is itself a theorem, proved elsewhere in the NEMS programme, but
importing that proof's premises (a formal notion of ``selection procedure invariant under
presentation'') into this repository's vocabulary is future formalization work, not a
philosophical shortcut. We flag this explicitly so that no reader mistakes the paper's overall
confidence in Theorem~5 for uniform machine-checked confidence in both of its halves.

\subsection{Relation to the digital-physics debate}

This paper's stance toward Zuse, Fredkin, Wolfram, and 't~Hooft on one side, and an
unrestricted-continuum view on the other, is neutrality of a specific and strong kind: not ``we
do not know which is right,'' but ``neither can be internally proved right, and this is a
theorem, not an admission of ignorance.'' That is a stronger and different claim than declining
to take a side in the ordinary sense. A digital-physics advocate who points to a successful
match between a candidate discrete rule and observed physics, however striking the match, is
producing exactly the situation \cref{thm:lawbarrier} addresses: a candidate law that has
matched so far, with no internal procedure available to certify that the match will continue
forever or that it is not one of the countably many discrete laws that happens to agree with
observations up to the current prefix while a continuum, or a different discrete law, diverges
downstream. Symmetrically, a continuum advocate who treats general relativity's differentiable
manifold as established fact about the actual substrate, rather than as the best-performing
finitely-described theory to date, is claiming a certification that \cref{thm:lawbarrier} shows
is unavailable to them too. The paper's neutrality is therefore not agnosticism about physics;
it is a proof that the entire debate, as usually conducted (each side citing empirical success as
if it amounted to internal proof), is arguing for something structurally unprovable from the
inside, regardless of which side is factually correct.

\subsection{Honest limitations beyond \texorpdfstring{\cref{sec:whatnot}}{Section~6}}

\cref{sec:whatnot} states the theorems' technical limits; here we state programmatic ones.
First, if computable analysis is eventually added to Mathlib, the refinement sketched in
\cref{rem:modeling} and \cref{sec:whatnot} item~4 should be carried out as a genuine
formalization project, not merely asserted via Pour-El--Richards; until then, the v1 development's
reliance on an unrestricted classical continuum is a real, disclosed simplification, not a
hidden one. Second, the quantum extension sketched in \cref{sec:open} --- replacing
\texttt{World} with a finite-dimensional versus separable Hilbert space and a POVM readout ---
is currently only a sketch: \cref{thm:prefix} is expected to go through essentially unchanged,
but whether \cref{thm:lawbarrier} survives with a unitary \(\mathrm{step}\) is a genuinely open
question, since the reduction's diagonal construction currently exploits classical branching
(\texttt{if}-\texttt{then}-\texttt{else} on a halting check) that does not obviously transfer to
unitary evolution without decoherence or measurement entering the model. Third, axis~C
(objective randomness) is addressed directly: the glimmer corollary (\cref{sec:glimmer}) is
proved in this paper, with its definitions and derivation chain machine-checked, because
\cref{thm:mdl-uniqueness}'s conclusion is not correctly statable without it. One honest gap
remains, and we disclose it precisely rather than rounding it up to a claimed result: the
Born-martingale lemma itself --- that the actual physical Born measure admits no successful
computable martingale on the realized transcript --- is not yet proved anywhere in the
programme. That is the \emph{only} gap: the Lean module proving the classification's
mathematical content (\leanref{Substrate/Glimmer.lean}) is part of this repository's audited,
sorry-free build, as disclosed in full in \cref{sec:glimmer}.

\subsection{Forward-looking directions}

Each remaining item in \cref{sec:open} matters for a distinct reason, stated there in full; in
summary: the deterministic half of the probabilistic-observers question is closed
(\cref{sec:open}), and Baire space carries a genuine probability measure under which
discreteness is measure-zero-typical (\cref{cor:measure-typicality}); what remains is genuinely
closer to real experimental practice --- a \emph{randomized} test with some advantage short of
certainty, evaluated against a measure that does not already assign full mass to one class by
construction. The cardinality strengthening of Theorem~5, \cref{cor:cardinality}, forecloses a
natural ``pool enough observers'' escape attempt, showing the barrier survives even an
uncountable family of cooperating internal observers. And the glimmer corollary's missing
Born-martingale lemma is the single result standing between this paper's programme and a
genuine resolution of axis~C's relationship to axis~B, which is why we name it explicitly as
the actual formal target rather than leaving axis~C as an unstructured gap.

% ============================================================================
% 10. OPEN PROBLEMS AND EXTENSIONS
% ============================================================================
\section{Open Problems and Extensions}
\label{sec:open}

\begin{itemize}[leftmargin=1.8em,itemsep=0.5em]
\item \textbf{Computable-analysis refinement of \texttt{ContinuumWorld}.} Restrict
\(\mathrm{step}\) to Cauchy-modulus-computable maps on Cauchy-modulus-computable reals.
\(\mathtt{Cstrict}\) provably remains nonempty under this refinement, by Pour-El--Richards
\citep{PourElRichards1981} (\cref{sec:history}, with the regularity caveat noted there).
Precisely what is missing: current Mathlib provides no computable-analysis framework at all
--- no notion of a Cauchy-modulus-computable real, no Weihrauch-degree machinery, nothing
analogous to the \texttt{Computable} predicate on \(\N \to \N\) already used throughout this
development (verified by direct search of the Mathlib source: zero occurrences of
\texttt{Weihrauch} or \texttt{CauchyModulus} anywhere in the library as of this writing).
Formalizing this refinement therefore means building a computable-analysis layer largely from
scratch, not applying existing machinery; that is well beyond this paper's scope.
\item \textbf{Quantum version.} Replace \texttt{World} with a finite-dimensional versus
separable Hilbert space and a POVM readout. \cref{thm:prefix} is expected to go through with
essentially the same proof strategy. Whether \cref{thm:lawbarrier} survives with a unitary
\(\mathrm{step}\) is open (\cref{sec:discussion}). Precisely what is missing: Mathlib currently
has zero formalization of POVMs, density operators, or completely-positive trace-preserving
(CPTP) maps (verified by direct search), so even \emph{stating} the quantum extension formally
--- not merely proving \cref{thm:lawbarrier}'s quantum analogue --- is future infrastructure
work, not a proof gap in the existing development.
\item \textbf{The Born-martingale lemma} (\cref{sec:glimmer}) --- the one genuinely open piece
of the glimmer corollary, whose definitions, four-class taxonomy, and derivation chain are
settled Part~II content (\cref{sec:glimmer}), not an open problem. What remains is a single
physical/interpretive premise: that the actual physical Born measure, as complete internal
semantics per NEMS Papers~13/14 \citep{SpivackP13,SpivackP14}, admits no successful computable
martingale on the realized transcript. This is not a randomness-theory gap --- the martingale
and computable-randomness machinery the lemma needs is already fully built and machine-checked
(\cref{sec:glimmer}) --- and it cannot be closed by definitional fiat without smuggling the
conclusion; see \cref{sec:glimmer} for the full derivation chain and precisely why this premise,
and only this premise, remains open.
\item \textbf{Probabilistic observers.} The deterministic class-advantage obstruction is
proved: an internal test has positive acceptance on \(\D\) if and only if it has positive
acceptance on \(\mathtt{Cstrict}\) (\leanref{Substrate.acceptance\_on\_D\_iff\_Cstrict}, module
\leanref{Substrate/Probabilistic.lean}), so no internal test exhibits exclusive class advantage
(\leanref{Substrate.not\_exclusive\_acceptance}) --- the prefix-lifting of \cref{thm:notfindet}
from properties to a test's acceptance behavior, by the identical mechanism
(\cref{thm:prefix}). Baire space carries a genuine measure-theoretic treatment:
\leanref{Substrate/Probabilistic.lean} wires in Mathlib's Ionescu--Tulcea-based
\texttt{MeasureTheory.Measure.infinitePi} directly
onto \(\N \to \N\), giving a canonical geometric i.i.d.\ product measure under which \(\D\) is
null and \(\mathtt{Cstrict}\) has full measure (\cref{cor:measure-typicality}). What remains
genuinely open is a \emph{genuinely randomized}, error-tolerant test. Prefix indistinguishability
(\cref{thm:prefix}) transports individual witnesses between classes, but it does not bound
\emph{aggregate} class-conditional acceptance-rate differences under a sampling measure: the
distribution of prefixes under a \(\D\)-conditioned sample and a \(\mathtt{Cstrict}\)-conditioned
sample can differ arbitrarily, so witness-transport alone gives no rate bound. A non-trivial
\(\varepsilon\)-tolerant result requires a genuinely randomized test --- one with independent
internal randomness, not merely a prefix-determined confidence function --- evaluated against a
\emph{mixture} measure that assigns controllable positive mass to both \(\D\) and
\(\mathtt{Cstrict}\) separately, since under any single i.i.d.\ product measure on the full
transcript space \(\D\) is null and \(\mathtt{Cstrict}\) has full mass
(\cref{cor:measure-typicality}), so conditioning on \(\D\) alone is degenerate under that
measure. A witness-based restatement using a prefix-determined confidence function
(\leanref{Substrate.HasProbAcceptAbove}) turns out to add no independent content on inspection:
since the confidence value depends only on the prefix, \cref{thm:prefix} transports
\emph{identical} confidence values between classes, not merely \(\varepsilon\)-close ones, so
the tolerance parameter never does any real work in that construction --- it is a notational
restatement of \leanref{Substrate.acceptance\_on\_D\_iff\_Cstrict}, not an independent
strengthening, and is omitted from Appendix~\ref{app:lean} on those grounds.
\end{itemize}

% ============================================================================
% 11. CONCLUSION
% ============================================================================
\section{Conclusion}
\label{sec:conclusion}

We have shown, and machine-checked in Lean~4 against Mathlib with zero \texttt{sorry} and zero
custom axioms, that no finite-resource observer internal to a physical system can decide
whether that system's substrate is discrete or continuum, on either of the two readings that
matter (information density and dynamics). The obstruction comes in two independent forms: a
topological barrier that needs no computability theory at all, and a recursion-theoretic barrier
that reduces from the halting problem even when the observer is handed a complete candidate law.
Under Perfect Self-Containment, this makes substrate discreteness an outsourced property, and the
only lawful internal-consistent adjudicator --- by an argument cited from the companion NEMS
programme, not re-derived here (\cref{sec:discussion}) --- is minimum-description-length
selection over complete law-sets: not a methodological preference, but the unique procedure that
survives the closure requirement. This dissolves the traditional grid-search experimental programme's
claim to settle the question, redirects empirical attention to asymptotic scaling exponents, and
places the epistemic status of finite, compressive physical theories --- including those coming
out of the UGP programme --- exactly where Theorem~5 puts it: strong evidence, never
certification. The results are not silent on objective randomness: the glimmer
corollary's definitions and separations are machine-checked (\cref{sec:glimmer}) --- but the
Born-martingale lemma's physical completeness premise, and the question of which substrate
actually obtains, are left, honestly, as open.

% ============================================================================
% APPENDIX A
% ============================================================================
\appendix

\section{Lean Statement List}
\label{app:lean}

\begin{quote}
\textbf{Axiom-printout status: final.} \texttt{Substrate/*.lean} builds cleanly
(\texttt{lake build}: 8262/8262 jobs, zero errors) and \texttt{grep -r sorry Substrate/} finds
nothing. The \texttt{\#print axioms} column below was obtained by running
\texttt{lake env lean scripts/axiom\_audit.lean} (the declarations that script covers include every \leanref{Substrate/Glimmer.lean} and \leanref{Substrate/AxisA.lean}
declaration named in \cref{sec:glimmer,rem:axisA}, plus \leanref{Substrate.law\_match\_pi2\_hard}
and \leanref{Substrate.law\_match\_pi2\_hard\_zero}) and, for the remaining declarations in this
table, an equivalent direct \texttt{\#print axioms} invocation against the built library, both
re-run independently for this draft. Every named theorem depends on at most the three standard
Mathlib axioms (\texttt{propext}, \texttt{Classical.choice}, \texttt{Quot.sound});
one (\leanref{Substrate.finitely\_determined\_compl}) depends on \emph{no} axioms at all. Zero
custom axioms, confirmed.
\end{quote}

\footnotesize
\setlength{\tabcolsep}{3pt}
\begin{longtable}{@{}>{\raggedright\arraybackslash}p{4.3cm}>{\raggedright\arraybackslash}p{6.7cm}>{\raggedright\arraybackslash}p{2.9cm}@{}}
\toprule
\textbf{Lean declaration} & \textbf{Statement} & \textbf{\#print axioms} \\
\midrule
\endfirsthead
\toprule
\textbf{Lean declaration} & \textbf{Statement} & \textbf{\#print axioms} \\
\midrule
\endhead
\midrule
\multicolumn{3}{r}{\textit{continued on next page}} \\
\endfoot
\bottomrule
\endlastfoot
\leanref{Substrate.transcript\_characterisation} & Lemma~2.1: \(f \in \D \leftrightarrow \mathrm{Computable}\ f\) & propext, Classical.choice, Quot.sound \\
\leanref{Substrate.discrete\_subset\_continuum} & Lemma~2.2: \(\D \subseteq \mathtt{C}\) & propext, Classical.choice, Quot.sound \\
\leanref{Substrate.mem\_C} & Every \(f : \N \to \N\) lies in \(\mathtt{C}\) & propext, Classical.choice, Quot.sound \\
\leanref{Substrate.exists\_noncomputable\_nat\_fn} & \(\exists f, \neg\,\mathrm{Computable}\ f\) & propext, Classical.choice, Quot.sound \\
\leanref{Substrate.not\_computable\_haltingIndicator} & Halting-indicator witness is not computable & propext, Classical.choice, Quot.sound \\
\leanref{Substrate.prefix\_indistinguishable} & Theorem~1: every prefix extends into both \(\D\) and \(\mathtt{Cstrict}\) & propext, Classical.choice, Quot.sound \\
\leanref{Substrate.D\_disjoint\_Cstrict} & \(\D\) and \(\mathtt{Cstrict}\) disjoint & propext, Classical.choice, Quot.sound \\
\leanref{Substrate.not\_finitely\_determined\_of\_dense\_pair} & Abstract dense-pair form of Theorem~2, for arbitrary disjoint prefix-dense \(A,B \subseteq \N \to \N\); no \texttt{Computable}/\texttt{Nat.Partrec} dependence in the proof term & propext, Classical.choice, Quot.sound \\
\leanref{Substrate.substrate\_not\_finitely\_determined} & Theorem~2: no finitely determined property separates \(\D\) from \(\mathtt{Cstrict}\), as the \(\D\)/\(\mathtt{Cstrict}\) instance of the abstract lemma above & propext, Classical.choice, Quot.sound \\
\leanref{Substrate.sound\_for\_continuum\_never\_fires} & Theorem~3(a): continuum-sound test never fires & propext, Classical.choice, Quot.sound \\
\leanref{Substrate.sound\_for\_discreteness\_never\_fires} & Theorem~3(b): discreteness-sound test never fires & propext, Classical.choice, Quot.sound \\
\leanref{Substrate.certifier\_silent\_or\_unsound} & Theorem~3 combined & propext, Classical.choice, Quot.sound \\
\leanref{Substrate.law\_match\_undecidable} & Theorem~4: law-match predicate is not computable, for any computable \(f\) & propext, Classical.choice, Quot.sound \\
\leanref{Substrate.law\_match\_undecidable\_zero} & Corollary~4.1: Theorem~4 specialized to \(f \equiv 0\) (Paper~30 bridge) & propext, Classical.choice, Quot.sound \\
\leanref{Substrate.law\_match\_pi2\_hard} & Theorem~4\('\): \(\mathrm{Tot} \le_m \mathrm{LawMatches}(f,\cdot)\) for any computable \(f\), via a sequential-composition reduction (run the candidate to completion, discard the result, emit \(f\)); \(\mathrm{Tot}\)'s \(\Pi^0_2\)-completeness is external, not proved here & propext, Classical.choice, Quot.sound \\
\leanref{Substrate.law\_match\_pi2\_hard\_zero} & Theorem~4\('\) specialized to \(f \equiv 0\) & propext, Classical.choice, Quot.sound \\
\leanref{Substrate.internally\_decidable\_implies\_finitely\_determined} & Internally decidable \(\Rightarrow\) finitely determined & propext, Classical.choice, Quot.sound \\
\leanref{Substrate.finitely\_determined\_compl} & Complement of a finitely determined property is finitely determined & \textbf{none} \\
\leanref{Substrate.substrate\_not\_internally\_decidable} & Theorem~5 (Lean half): no internally decidable property separates \(\D\) from \(\mathtt{Cstrict}\) & propext, Classical.choice, Quot.sound \\
\leanref{Substrate.substrate\_not\_internally\_decidable'} & Symmetric separation & propext, Classical.choice, Quot.sound \\
\leanref{Substrate.discrete\_not\_internally\_decidable} & \(f \in \D\) not internally decidable & propext, Classical.choice, Quot.sound \\
\leanref{Substrate.strict\_continuum\_not\_internally\_decidable} & \(f \in \mathtt{Cstrict}\) not internally decidable & propext, Classical.choice, Quot.sound \\
\leanref{Substrate.D\_countable} & \(\D\) is countable & propext, Classical.choice, Quot.sound \\
\leanref{Substrate.Cstrict\_not\_countable} & \(\mathtt{Cstrict}\) is not countable & propext, Classical.choice, Quot.sound \\
\leanref{Substrate.no\_joint\_substrate\_decision} & \cref{cor:cardinality}: no family of internal tests, over any index type, jointly decides discrete membership & propext, Classical.choice, Quot.sound \\
\leanref{Substrate.acceptance\_on\_D\_iff\_Cstrict} & Deterministic class-advantage obstruction: positive acceptance on \(\D\) iff positive acceptance on \(\mathtt{Cstrict}\) & propext, Classical.choice, Quot.sound \\
\leanref{Substrate.D\_null\_geometric} & \cref{cor:measure-typicality}: \(\D\) has measure zero under the geometric i.i.d.\ product on Baire space & propext, Classical.choice, Quot.sound \\
\leanref{Substrate.Cstrict\_full\_mass\_geometric} & \cref{cor:measure-typicality}: \(\mathtt{Cstrict}\) has full measure under the same measure & propext, Classical.choice, Quot.sound \\
\leanref{Substrate.limitComputable\_not\_finitely\_determined} & \cref{sec:glimmer}: the abstract dense-pair form of Theorem~2 applied to limit-computable vs.\ not-limit-computable sequences & propext, Classical.choice, Quot.sound \\
\leanref{Substrate.computable\_not\_random} & \cref{sec:glimmer}, derivation-chain step~2's supporting fact: a computable, non-degenerate \(f\) admits a computable \(\mu\)-martingale that succeeds, hence is not computably random & propext, Classical.choice, Quot.sound \\
\leanref{Substrate.glimmer\_not\_computable} & \cref{sec:glimmer}: a glimmer sequence is not computable & propext, Classical.choice, Quot.sound \\
\leanref{Substrate.glimmer\_disjoint\_D} & \cref{sec:glimmer}: \(\mathrm{Glimmer}\) is disjoint from \(\D\) & propext, Classical.choice, Quot.sound \\
\leanref{Substrate.glimmer\_subset\_Cstrict} & \cref{sec:glimmer}: every glimmer sequence lies in \(\mathtt{Cstrict}\) & propext, Classical.choice, Quot.sound \\
\leanref{Substrate.glimmer\_finite\_description} & \cref{sec:glimmer}: a limit-computable sequence is fixed by one code plus a convergent stage function (non-vacuous finite-description witness) & propext, Classical.choice, Quot.sound \\
\leanref{Substrate.BetMartingale.bet\_succeeds} & \cref{sec:glimmer}: the bet-everything martingale construction succeeds on any computable, non-degenerate \(f\) & propext, Classical.choice, Quot.sound \\
\leanref{Substrate.Dfin\_iff\_eventuallyPeriodic} & \cref{rem:axisA}: \(\mathtt{Dfin}\) coincides with the eventually periodic sequences & propext, Classical.choice, Quot.sound \\
\leanref{Substrate.axisA\_not\_finitely\_determined} & \cref{rem:axisA}: the abstract dense-pair form of Theorem~2 applied to \(\mathtt{Dfin}\) vs.\ its complement & propext, Classical.choice, Quot.sound \\
\end{longtable}
\normalsize

\noindent\textbf{Module map.}

\footnotesize
\begin{longtable}{@{}>{\raggedright\arraybackslash}p{4.0cm}>{\raggedright\arraybackslash}p{9.9cm}@{}}
\toprule
\textbf{Module} & \textbf{Role} \\
\midrule
\endfirsthead
\toprule
\textbf{Module} & \textbf{Role} \\
\midrule
\endhead
\midrule
\multicolumn{2}{r}{\textit{continued on next page}} \\
\endfoot
\bottomrule
\endlastfoot
\leanref{Substrate/Basic.lean} & \leanref{Substrate.World}, transcript, \leanref{Substrate.D} / \leanref{Substrate.C} / \leanref{Substrate.Cstrict}, \leanref{Substrate.InternalTest}, \leanref{Substrate.FinitelyDetermined} \\
\leanref{Substrate/Discrete.lean} & Lemma~2.1 \\
\leanref{Substrate/Embed.lean} & Lemma~2.2 \\
\leanref{Substrate/Noncomputable.lean} & Existence of a non-computable \(\N \to \N\) function, via the halting problem \\
\leanref{Substrate/Prefix.lean} & Theorem~1 and the disjointness/density corollary \\
\leanref{Substrate/NotFinDet.lean} & Theorem~2, and its abstract dense-pair generalization (\cref{thm:notfindet}'s keybox) \\
\leanref{Substrate/Certifier.lean} & Theorem~3 \\
\leanref{Substrate/LawBarrier.lean} & Theorem~4 (halting reduction via \texttt{evaln}) and Theorem~4\('\) (\(\Pi^0_2\)-hardness via the sequential-composition reduction) \\
\leanref{Substrate/Outsourced.lean} & Theorem~5 (Lean-checkable half) \\
\leanref{Substrate/Cardinality.lean} & \cref{cor:cardinality}: countability of \(\D\), non-countability of \(\mathtt{Cstrict}\), and the uncountable-observer strengthening \\
\leanref{Substrate/Probabilistic.lean} & Deterministic class-advantage obstruction for probabilistic observers, and measure-zero / full-measure classification of \(\D\) / \(\mathtt{Cstrict}\) under a geometric product measure (\cref{cor:measure-typicality,sec:open}) \\
\leanref{Substrate/Glimmer.lean} & \cref{sec:glimmer}: limit-computability, computable Born measures, computable martingales, computable randomness, and the glimmer class, with the bet-everything martingale construction and the class-1/2\('\)/2/3 separation theorems \\
\leanref{Substrate/AxisA.lean} & \cref{rem:axisA}: the axis-A dense pair \(\mathtt{Dfin}\) (eventually periodic sequences) and its own instance of \cref{thm:notfindet} \\
\end{longtable}
\normalsize

\section{Epistemic Status Table}
\label{app:status}

\footnotesize
\begin{longtable}{@{}>{\raggedright\arraybackslash}p{9.5cm}>{\raggedright\arraybackslash}p{4.4cm}@{}}
\toprule
\textbf{Claim} & \textbf{Status} \\
\midrule
\endfirsthead
\toprule
\textbf{Claim} & \textbf{Status} \\
\midrule
\endhead
\midrule
\multicolumn{2}{r}{\textit{continued on next page}} \\
\endfoot
\bottomrule
\endlastfoot
Lemma~2.1, Lemma~2.2; Theorems~1--4; Corollary~1.1, Corollary~4.1 & \A~(machine-checked, Lean~4 / Mathlib, zero \texttt{sorry}, zero custom axioms) \\
Abstract dense-pair form of Theorem~2 (\cref{thm:notfindet}'s keybox) & \A~(module \leanref{Substrate/NotFinDet.lean}; no \texttt{Computable}/\texttt{Nat.Partrec} dependence in the proof term) \\
Theorem~4\('\) (\(\Pi^0_2\)-hardness of \(\mathrm{LawMatches}\), \leanref{Substrate.law\_match\_pi2\_hard}/\leanref{Substrate.law\_match\_pi2\_hard\_zero}) & \A~(module \leanref{Substrate/LawBarrier.lean}; \(\mathrm{Tot} \le_m \mathrm{LawMatches}(f,\cdot)\) via a sequential-composition reduction; \(\mathrm{Tot}\)'s \(\Pi^0_2\)-completeness is an external, cited recursion-theoretic fact, not proved in this repository) \\
Theorem~5, outsourcing half (\cref{thm:outsourcing}) & \A \\
Theorem~5, MDL-uniqueness half (\cref{thm:mdl-uniqueness}) & \AMDL~(argument cited to NEMS Papers~34/84; not a Mathlib reduction) \\
Uncountable-observer strengthening of Theorem~5 (\cref{cor:cardinality}) & \A~(module \leanref{Substrate/Cardinality.lean}; closed no bounds on the index type \(I\), including uncountable) \\
Probabilistic observers, deterministic class-advantage obstruction (\leanref{Substrate.acceptance\_on\_D\_iff\_Cstrict}) & \A~(module \leanref{Substrate/Probabilistic.lean}) \\
Probabilistic observers, measure-zero / full-measure classification of \(\D\) / \(\mathtt{Cstrict}\) under a geometric product measure (\cref{cor:measure-typicality}) & \A~(module \leanref{Substrate/Probabilistic.lean}; uses Mathlib's Ionescu--Tulcea-based \texttt{Measure.infinitePi}) \\
Probabilistic observers, \(\varepsilon\)-tolerant \emph{randomized} generalization (genuine internal randomness plus a class-charged mixture measure) & \textbf{Open} --- prefix indistinguishability transports individual witnesses but not aggregate class-conditional acceptance rates; a prefix-determined witness restatement (\leanref{Substrate.HasProbAcceptAbove}) adds no independent content and is omitted from Appendix~\ref{app:lean} on those grounds (\cref{sec:open}) \\
\S\ref{sec:scaling} scaling-exponent classification & \textsc{d}~(discussion-level, not machine-checked, presented with citations) \\
\S\ref{sec:glimmer} glimmer definitions (\leanref{Substrate.LimitComputable}, \leanref{Substrate.BornMeasure}, \leanref{Substrate.Glimmer}, etc.), limit-computability density pair, \leanref{Substrate.limitComputable\_not\_finitely\_determined}, \leanref{Substrate.glimmer\_finite\_description}, \leanref{Substrate.computable\_not\_random}, the bet-everything martingale construction (\leanref{Substrate.BetMartingale.bet}/\leanref{Substrate.BetMartingale.bet\_succeeds}), and the separations built on it (\leanref{Substrate.glimmer\_not\_computable}, \leanref{Substrate.glimmer\_disjoint\_D}, \leanref{Substrate.glimmer\_subset\_Cstrict}) & \A~(module \leanref{Substrate/Glimmer.lean}, part of the audited, sorry-free build; imported by \leanref{Substrate.lean}; every declaration confirmed by \texttt{\#print axioms} to depend on only \texttt{propext}, \texttt{Classical.choice}, \texttt{Quot.sound}) \\
\S\ref{sec:glimmer} Glimmer non-emptiness (Chaitin's \(\Omega\)) & \A/\textsc{d}~(cited to Downey--Hirschfeldt \citep{DowneyHirschfeldt2010}, not built, by design) \\
\S\ref{sec:glimmer} Born-martingale lemma (the physical completeness premise) & \textbf{Open} --- not proved here or anywhere in the programme at time of writing. Not a randomness-theory gap: the martingale/computable-randomness machinery (\leanref{Substrate.computable\_not\_random}, \leanref{Substrate.BetMartingale.bet\_succeeds}) is machine-checked (row above); the unrestricted universal form of the premise is in fact \emph{false} here (\leanref{Substrate.BetMartingale.bet\_succeeds} beats every \emph{computable} non-degenerate \(f\)), so defining ``complete internal semantics'' as ``\(f\) is computably random'' would smuggle the conclusion. The nearest existing candidate, \texttt{nems-lean}'s \texttt{BICS} (Papers~13/14), is a quantum semantic-architecture predicate over density operators, categorically the wrong kind of object for a realized-sequence bridge. What is missing is specifically a non-question-begging formalization of the physical premise as a predicate on realized transcripts (\cref{sec:open}) \\
\cref{rem:axisA} axis-A pair (\leanref{Substrate.Dfin}, \leanref{Substrate.Dfin\_iff\_eventuallyPeriodic}, \leanref{Substrate.axisA\_not\_finitely\_determined}) & \A~(module \leanref{Substrate/AxisA.lean}, part of the audited, sorry-free build; imported by \leanref{Substrate.lean}; every declaration confirmed by \texttt{\#print axioms} to depend on only \texttt{propext}, \texttt{Classical.choice}, \texttt{Quot.sound}) \\
\end{longtable}
\normalsize

% ============================================================================
% BIBLIOGRAPHY
% ============================================================================
\bibliographystyle{plainnat}
\bibliography{refs}

\end{document}
